%
% messaging-specification.tex
% Compile: pdflatex messaging-specification.tex
%

\documentclass[a4paper,10pt,fleqn]{article}
\usepackage[margin=1in]{geometry}
\usepackage[colorlinks=true,linkcolor=blue,citecolor=blue,urlcolor=blue]{hyperref}
\usepackage{booktabs}
\usepackage{array}
\usepackage{tabularx}
\usepackage{longtable}
\usepackage{enumitem}
\usepackage{listings}
\usepackage{xcolor}
\usepackage{amsmath}
\usepackage{microtype}

\setlength{\parskip}{0.7\baselineskip plus 0.2\baselineskip minus 0.1\baselineskip}
\setlength{\parindent}{0pt}

\newcolumntype{L}[1]{>{\raggedright\arraybackslash}p{#1}}
\newcolumntype{Y}{>{\raggedright\arraybackslash}X}

\lstset{
  basicstyle=\ttfamily\small,
  keywordstyle=\color{blue!80!black},
  commentstyle=\color{green!50!black},
  stringstyle=\color{red!60!black},
  breaklines=true,
  frame=single,
  framerule=0.3pt,
  rulecolor=\color{gray!50},
  backgroundcolor=\color{gray!5},
  columns=fullflexible,
}

\newcommand{\cls}[1]{\texttt{#1}}
\newcommand{\sel}[1]{\texttt{\##1}}
\newcommand{\pkg}[1]{\texttt{#1}}
\newcommand{\goref}[1]{\textsf{\small [Go: #1]}}

\begin{document}

\title{Claude Messaging SDK for Pharo:\\Technical Specification}
\author{Punt Labs}
\date{April 2026 --- v0.3.0-dev}
\hypersetup{
  pdftitle={Claude Messaging SDK for Pharo: Technical Specification},
  pdfauthor={Punt Labs},
  pdfsubject={Messaging SDK Specification},
  pdfcreator={pdflatex}
}
\maketitle

\tableofcontents
\newpage

\subsection*{About this Specification}

Section order mirrors Anthropic's Messages API docs
(\url{https://docs.claude.com/en/api/messages}).  Unimplemented
surface is marked \textbf{Gap (Phase 2)}; see
\texttt{messaging-specification-gap-analysis.md}.  Pharo-specific
implementation notes live in
\emph{messaging-specification-pharo-notes.tex}.  The Agent layer
(tool runner, sub-agents, dashboard) is in the separate
\emph{Claude Agent SDK for Pharo --- Technical Specification}.


\part{Learn about Claude}

\section{Overview}
\label{sec:overview}

Anthropic's Messages API is the primary endpoint for chat-style
generation against Claude models.  A caller sends an ordered list
of user and assistant turns; the API returns a typed response
carrying assistant text, tool-use blocks, citations, stop reason,
and token usage.  Streaming is a transport variant of the same
shape, delivering the response as Server-Sent Events (SSE).
Synchronous chat-style generation, tool-use loops, and streaming
interactions route through this endpoint; asynchronous batch
processing reuses the Messages request shape but submits via the
Batches API instead.  The other endpoint families --- Files,
Skills, Admin, Batches, Prompt Tools, Models, and Count Tokens
--- live as separate surfaces that compose with Messages or
stand alone.

The API is built for the workloads people actually run against
Claude: agents that loop tool calls back into the model, copilots
that assist a developer in an editor, chat UIs that hold a
turn-based conversation, retrieval-augmented generation pipelines
that ground answers in fetched documents, document Q\&A,
classification, summarization, structured extraction, and content
generation.  It is also the backbone for the rest of Anthropic's
platform: server-side tools attach via the \texttt{tools} array
in the request body (no beta), MCP connectors attach via both the
\texttt{mcp\_servers} request body field and the
\texttt{mcp-client-2025-04-04} beta header, and Files and Skills
attach via their own beta headers (\texttt{files-api-2025-04-14},
\texttt{skills-2025-10-02}) that opt the request into a feature group.

The Claude Messaging SDK for Pharo is a Smalltalk client for
that API, running inside a live Pharo~12 image.  It
gives a Smalltalker the same surface a Python or TypeScript developer
gets from the first-party SDKs --- build a message, send it, read
Claude's response as a typed object --- while keeping to idiomatic
Smalltalk conventions throughout: fluid class syntax, typed hierarchies
instead of dictionaries, \texttt{>>>} doctest examples, Iceberg-managed
Tonel packages.  The SDK ships in two layers.  This document specifies
the \emph{Messaging layer}: \texttt{ClaudeClient}, the
\texttt{ClaudeMessage} / \texttt{ClaudeContentBlock} hierarchy, the
error classes, the streaming decoder, the beta-gated APIs (Files,
Skills, MCP).  The \emph{Agent layer} --- tool runner, sub-agents,
dashboard, rate limiter --- builds on this surface and is described in
the separate \emph{Claude Agent SDK for Pharo} specification.

The SDK covers the public Messages API endpoints at
\texttt{https://api.anthropic.com} (configurable via \texttt{baseUrl}
on \texttt{ClaudeClient} for proxies or mocks): \texttt{POST
/v1/messages} for creating messages (streaming and non-streaming),
\texttt{POST /v1/messages/count\_tokens} for token counting,
\texttt{GET /v1/models} and \texttt{/v1/models/\{id\}} for model
discovery, and the Files API endpoints (\texttt{/v1/files} for upload,
list, get, download, and delete).  The endpoint map in full:

\begin{center}
\begin{tabular}{@{}lll@{}}
\toprule
\textbf{Endpoint} & \textbf{Method} & \textbf{Description} \\
\midrule
\texttt{/v1/messages}              & POST & Create message (streaming and non-streaming) \\
\texttt{/v1/messages/count\_tokens} & POST & Count tokens for a message \\
\texttt{/v1/models}                & GET  & List available models \\
\texttt{/v1/models/\{id\}}        & GET  & Get model details \\
\texttt{/v1/files}               & POST & Upload a file \\
\texttt{/v1/files}               & GET  & List uploaded files \\
\texttt{/v1/files/\{id\}}        & GET  & Get file metadata \\
\texttt{/v1/files/\{id\}/content} & GET & Download file content \\
\texttt{/v1/files/\{id\}}        & DELETE & Delete a file \\
\bottomrule
\end{tabular}
\end{center}

The example below is the canonical first SDK call: build a client with
an API key, build a request with a model and a user message, send it,
and read the text off the response.  Every later section elaborates on
one piece of this shape --- request options, streaming, tool use,
prompt caching, thinking --- but a reader who gets this example running
has the whole SDK in hand.  The example evaluates to Claude's reply
text as a String:

\begin{lstlisting}[language={}]
| client request response |
"Build a client (see ClaudeSDKExampleSupport resolveClient for keyring
 resolution in scripts; the apiKey: factory is the minimal shape)."
client := ClaudeClient apiKey: 'sk-ant-...'.

"Build a request: model, max tokens, one user turn."
request := ClaudeMessageRequest new
    model: ClaudeModel opus46;
    maxTokens: 1024;
    addUserMessage: 'Write a haiku about Smalltalk.';
    yourself.

"Synchronous POST to /v1/messages; returns a ClaudeMessage."
response := client sendMessage: request.

"textContent joins all text blocks into a single String."
response textContent
\end{lstlisting}

Pharo-specific implementation notes --- threading discipline around
the Morphic UI, Iceberg / Tonel packaging, image save and restore
behavior, supervised process patterns --- live in the companion
document \emph{Claude Messaging SDK for Pharo: Pharo-Specific
Implementation Notes}.

For the full shape of the Messages API surface see
\S\ref{sec:messages}; for streaming see \S\ref{sec:streaming}; for
thinking see \S\ref{sec:thinking}.

\medskip
\noindent{\small\textsf{Status: \emph{Available}.\quad
Anthropic docs: \url{https://docs.claude.com/en/api/messages}.}}

\section{Design Goals}
\label{sec:design-goals}

The SDK is shaped by four commitments that trade off against each
other in predictable ways: parity with the Messages API defines the
surface, idiomatic Smalltalk defines the style, Pharo-safe threading
defines the concurrency contract, and a minimal dependency footprint
keeps the package graph light enough to load into any Pharo~12 image.
The numbered list below names each commitment and the specific choices
that follow from it.

\begin{enumerate}[label=\textbf{G\arabic*}.]
  \item \textbf{Feature parity with Anthropic's Messages API.}  Every
    v1.0 Pharo-scope feature in \S\ref{sec:scope} maps to a typed
    Smalltalk surface: the Messages API itself (send, stream, tool
    use), the streaming decoder, models API, token counting, the error
    hierarchy, core options, extended thinking, citations, server tool
    results, server-side tools, the MCP connector, Skills API, Files
    API, output encoding, and the message-decomposition primitives
    (turn, stop details, request context, container).  Both stable and
    beta APIs ship together --- beta headers are managed per
    \S\ref{sec:beta-headers}, automatic for Files and Skills, opt-in
    via \sel{betas:} for others --- so a caller can reach the full
    v1.0 surface without conditional code paths.
  \item \textbf{Idiomatic Smalltalk.}  Follow Kent Beck's
    \emph{Smalltalk Best Practice Patterns}: composed methods,
    intention-revealing selectors, pluggable behavior, collecting
    parameters.  Use the class hierarchy to carry type information
    rather than dictionaries or type tags --- \texttt{ClaudeMessage}
    and \texttt{ClaudeContentBlock} are families with polymorphic
    dispatch, not schema-bag records.  JSON serialization is explicit
    (\texttt{asJson} / \texttt{fromJson:} with a \texttt{jsonKeyMap})
    rather than reflective.  The SDK should read like Smalltalk to a
    Smalltalker.
  \item \textbf{Pharo-safe threading.}  Network I/O must never block
    the Morphic UI thread.  The SDK itself is synchronous, so every
    caller knows exactly which selectors block and for how long; all
    blocking operations are documented so consumers know what to fork
    with \texttt{forkAt:} and what to defer back to the UI via
    \texttt{WorldState defer:}.  The companion Pharo-notes document
    covers the forking patterns.
  \item \textbf{Minimal dependencies.}  Use only Pharo's built-in
    libraries: Zinc for HTTP, STONJSON for JSON, and the standard
    collections.  No external Metacello baselines, no FFI to OS
    libraries from the Messaging layer.  The SDK loads into a stock
    Pharo~12 image by filing in the \texttt{Claude-Messaging-*} Tonel
    packages from \texttt{src/} (driven by \texttt{make setup} /
    \texttt{make filein}, which use \texttt{TonelReader} and
    \texttt{MCPackageLoader}); per-product Metacello baselines are
    future work.
\end{enumerate}

This section is reference material; no runnable example applies.

\medskip
\noindent{\small\textsf{Status: \emph{Available}.}}

\section{Scope}
\label{sec:scope}

This SDK covers the public Messages API plus the beta APIs that
stabilized as of v1.0.  The table below enumerates each feature with
its v1.0 Pharo scope and the version the Pharo surface landed in.
Features marked for a later version column ship later --- the SDK
reserves version slots so callers can track when a surface becomes
available rather than guessing from a changelog.  Features marked
\emph{Out of scope} belong to adjacent Anthropic surfaces (Bedrock,
Vertex, legacy completions) or to the Agent layer; the scope boundary
is deliberate and documented per row.

\begin{center}
\begin{tabularx}{\linewidth}{@{}lYl@{}}
\toprule
\textbf{Feature} & \textbf{Smalltalk Scope} & \textbf{Version} \\
\midrule
Messages API          & Full (send, stream, tool use)        & v1 \\
Streaming (SSE)       & Full (decoder, events, accumulator)  & v1 \\
Models API            & List, get, model info type           & v1 \\
Token counting        & CountTokens method and return type   & v1 \\
Error handling        & Full hierarchy, retry, rate limits   & v1 \\
Options system        & Core options (key, URL, timeout)     & v1 \\
Extended thinking     & ThinkingConfig, thinking blocks, redacted & v1 \\
Citations             & Citation types on TextBlock          & v1 \\
Server tool results   & Read-only response types             & v1 \\
Server-side tools     & Full (web search, web fetch, code exec, bash, text editor, memory, tool search) & v1 \\
MCP connector         & Full (request param + response content blocks) & v1 \\
Skills API            & Full (CRUD + versions + container loading + multipart upload) & v1 \\
Files API             & Full (upload, list, get, download, delete) & v1 \\
Output encoding       & Text output format (encoding + markup) & v1 \\
Message decomposition & Turn, stop details, request context, container & v1 \\
Tool runner           & See \emph{Claude Agent SDK Specification} & v1 \\
Message batches       & Not yet implemented                  & v2 \\
Bedrock / Vertex      & Out of scope                         & --- \\
Completions (legacy)  & Out of scope                         & --- \\
\bottomrule
\end{tabularx}
\end{center}

The Batches API row is tracked against the full specification in
\S\ref{sec:batches}; the Files API lands in \S\ref{sec:files}.

This section is reference material; no runnable example applies.

\medskip
\noindent{\small\textsf{Status: \emph{Available}.\quad
Anthropic docs: \url{https://docs.claude.com/en/api/messages}.}}

\part{API Reference}

\section{Messages API}
\label{sec:messages}

The Messages API is Claude's primary generation endpoint --- every
synchronous chat-style interaction, tool-use loop, and streaming
response rides on \texttt{POST /v1/messages}.  The Pharo SDK
exposes that endpoint through three collaborating objects:
\texttt{ClaudeClient} holds authentication and connection
configuration, \texttt{ClaudeMessageRequest} builds the request
body (model, messages, options), and \texttt{ClaudeMessage}
comes back as the typed response.  Any synchronous chat-style
generation, tool-use loop, or streaming interaction routes through
here; asynchronous batch processing reuses the same Messages request
shape but submits via the Batches API endpoints instead.  The other
endpoint families (Files, Skills, Admin, Batches, Prompt Tools,
Models, Count Tokens) live as separate HTTP endpoint surfaces that
compose with Messages or stand alone.  Thinking, tool use,
streaming, prompt caching, and citations are all features
layered on top of a Messages request.

A \texttt{ClaudeClient} is cheap to build and meant to be reused.
The simplest form takes an API key; the three-keyword factory
adds a base URL override (for proxies or mock servers) and a
timeout.  \texttt{fromEnvironment} reads the
\texttt{ANTHROPIC\_API\_KEY} variable via \texttt{Smalltalk os
environment} and raises if it is unset --- useful for scripts
and CI runners where the key is injected from outside the image.

\begin{lstlisting}[language={}]
"Explicit API key."
client := ClaudeClient apiKey: 'sk-ant-...'.

"From environment variable ANTHROPIC_API_KEY."
client := ClaudeClient fromEnvironment.

"With a proxy or mock server URL plus a tighter timeout."
client := ClaudeClient
    apiKey: 'sk-ant-...'
    baseUrl: 'https://custom-proxy.example.com'
    timeout: 120.
\end{lstlisting}

The client's configuration lives in eight instance variables, all
with defaults set by \texttt{initialize}:

\begin{center}
\small
\begin{tabular}{@{}llp{5.5cm}@{}}
\toprule
\textbf{Variable} & \textbf{Type} & \textbf{Default} \\
\midrule
\texttt{apiKey}    & String   & From \sel{apiKey:} argument or \texttt{fromEnvironment} \\
\texttt{baseUrl}   & String   & \texttt{https://api.anthropic.com} \\
\texttt{apiVersion} & String  & \texttt{2023-06-01} \\
\texttt{timeout}   & Integer  & 120 (seconds); used for non-streaming requests \\
\texttt{streamingTimeout} & Integer & 300 (seconds); used for \sel{streamMessage:do:}
  and as the socket read timeout in \sel{openStreamingSocket:} \\
\texttt{maxRetries} & Integer & 5 \\
\texttt{rateLimiter} & \cls{ClaudeRateLimiter} or nil & Optional
  shared rate limiter; populated when the client joins an Agent-layer
  rate limiter \\
\texttt{debugLogBlock} & Block or nil & Optional callable that
  receives a Dictionary per request and response event (keys
  \texttt{\#direction}, \texttt{\#url}, \texttt{\#body},
  \texttt{\#status}). Wire with \sel{debugLog:}. \\
\bottomrule
\end{tabular}
\end{center}

A fresh \texttt{ZnClient} is created for each request rather than
pooled, so every call performs its own TCP/TLS handshake.  That
is acceptable for an API whose typical conversation turn makes
one to five requests; it also means \texttt{ClaudeClient} survives
Pharo image save without special handling --- no connections are
held across save.  A process that was mid-request at save time
wakes up on a broken socket after restore; long-running consumers
should guard their network forks accordingly.

\subsection{Request Lifecycle}
\label{sec:request-lifecycle}

Every Messages request follows the same seven-step pipeline, and
understanding it is how you reason about retry, timeout, and
error shape.  The request object is serialized to JSON with
\texttt{STONJSON toString:}; a fresh \texttt{ZnClient} is
configured with \texttt{x-api-key}, \texttt{anthropic-version:
2023-06-01}, \texttt{content-type: application/json}, and
\texttt{accept: application/json}; the POST fires; any non-2xx
response is inspected for retry eligibility and either retried
with backoff (see \S\ref{sec:errors}) or raised as a typed
\texttt{ClaudeApiError} subclass; the response body is parsed
with \texttt{STONJSON fromString:}; the resulting dictionary is
handed to \texttt{ClaudeMessage class>>fromJson:}; and the typed
\texttt{ClaudeMessage} comes back to the caller.  The HTTP POST
and response-parse steps are blocking network I/O.  The SDK
itself never forks --- consumers choose the concurrency strategy
that fits their application, which means a call from the Morphic
UI thread must be wrapped in \texttt{forkAt:} to avoid locking
the UI for the duration of the request.

\subsection{Sending a Message}

The end-to-end example below covers the \texttt{sendMessage:}
path against a \texttt{ClaudeMockServer} so it runs without an
API key or network.  Real code swaps the mock for
\texttt{ClaudeSDKExampleSupport resolveClient} and removes the
\texttt{baseUrl:} override.  The request exercises three
delegated setters beyond \S\ref{sec:overview}'s minimal hello
call: a system prompt, a sampling temperature, and a metadata
tag that the server echoes in usage analytics.  The example
evaluates to an array with three fields pulled off the typed
response --- text, stop reason symbol, and input-token count:

\begin{lstlisting}[language={}]
| server client request response |
"Spin up an in-image mock HTTP server so the example needs neither network nor API key."
server := ClaudeMockServer startOn: 0.
[ "Register the canned response for POST /v1/messages."
  server
    respondTo: '/v1/messages'
    with: '{"id":"msg_01","type":"message","role":"assistant","content":[{"type":"text","text":"2 plus 2 is 4."}],"model":"claude-sonnet-4-5","stop_reason":"end_turn","stop_sequence":null,"usage":{"input_tokens":18,"output_tokens":7,"cache_creation_input_tokens":0,"cache_read_input_tokens":0}}'.
  "Build a ClaudeClient that talks to the mock."
  client := ClaudeClient
    apiKey: 'test-key'
    baseUrl: 'http://localhost:' , server port printString
    timeout: 5.
  "Build the request with a system prompt, a low temperature, and a metadata tag."
  request := ClaudeMessageRequest new
    model: ClaudeModel sonnet45;
    maxTokens: 256;
    system: 'You are a concise arithmetic tutor.';
    temperature: 0.0;
    metadata: (Dictionary newFromPairs: { 'user_id'. 'example-reader' });
    addUserMessage: 'What is 2 + 2?';
    yourself.
  "Synchronous POST to /v1/messages; returns a ClaudeMessage."
  response := client sendMessage: request.
  ^ { response textContent. response stopReason. response usage inputTokens } ]
    ensure: [ server stop ]  "Always stop the mock, even on error."
"=> #('2 plus 2 is 4.' #end_turn 18)"
\end{lstlisting}

Application code wraps the \texttt{sendMessage:} call in an
\texttt{on: ClaudeApiError do:} handler when it wants to react
to specific error subclasses; for most callers the default
behavior (raise the typed exception) is correct because
retryable errors are already retried inside the client (see
\S\ref{sec:errors}).

\subsection{Request Parameters (ClaudeMessageRequest)}

\cls{ClaudeMessageRequest} is the request body as a first-class
Smalltalk object.  It holds the two required fields directly
(\texttt{model}, \texttt{maxTokens}) along with the conversation
turns and system prompt, and delegates optional parameters to
five composite value objects.  Breaking the surface into
composites (sampling knobs, tool parameters, thinking params,
request options, output config) keeps the top-level class small
enough to scan while letting each family of options evolve on
its own schedule:

\begin{center}
\small
\begin{tabular}{@{}llp{6cm}@{}}
\toprule
\textbf{Variable} & \textbf{Type} & \textbf{Description} \\
\midrule
\texttt{model}          & String                     & Model ID (required) \\
\texttt{maxTokens}      & Integer                    & Max output tokens (required) \\
\texttt{messages}       & OrderedCollection          & Conversation turns \\
\texttt{system}         & String or nil              & System prompt \\
\texttt{sampling}       & \cls{ClaudeSamplingParams} or nil & Sampling knobs (temperature, topP, topK, stopSequences) \\
\texttt{toolParams}     & \cls{ClaudeToolParams} or nil     & Tool declarations and tool-choice policy \\
\texttt{thinkingParams} & \cls{ClaudeThinkingParams} or nil & Extended thinking mode (enabled, adaptive, disabled) \\
\texttt{requestOptions} & \cls{ClaudeRequestOptions} or nil & Per-request controls (metadata, cache control, service tier, betas, MCP servers, container, etc.) \\
\texttt{outputConfig}   & \cls{ClaudeOutputConfig} or nil   & Effort hint and output format (text or JSON schema) \\
\bottomrule
\end{tabular}
\end{center}

The composites are initialized in \texttt{initialize}, so a
bare \texttt{ClaudeMessageRequest new} already has empty
sampling, tool, thinking, request-options, and output-config
objects ready to accept setters.  Callers never need to
instantiate the composites themselves; the request forwards the
common setters directly, and each forwarded message sets the
field on the right composite:

\begin{center}
\small
\begin{tabular}{@{}p{4.2cm}lp{4.8cm}@{}}
\toprule
\textbf{Selectors} & \textbf{Held by} & \textbf{Wire field(s)} \\
\midrule
\sel{temperature:} \sel{topP:} \sel{topK:} \sel{stopSequences:}
   & \cls{ClaudeSamplingParams}
   & \texttt{temperature}, \texttt{top\_p}, \texttt{top\_k}, \texttt{stop\_sequences} \\
\sel{tools:} \sel{toolChoice:}
   & \cls{ClaudeToolParams}
   & \texttt{tools}, \texttt{tool\_choice} \\
\sel{thinking:}
   & \cls{ClaudeThinkingParams}
   & \texttt{thinking} \\
\sel{metadata:} \sel{cacheControl:} \sel{serviceTier:} \sel{inferenceGeo:} \sel{speed:} \sel{betas:} \sel{mcpServers:} \sel{container:} \sel{stream:}
   & \cls{ClaudeRequestOptions}
   & \texttt{metadata}, \texttt{cache\_control}, \texttt{service\_tier}, \texttt{inference\_geo}, \texttt{speed}, \texttt{anthropic-beta} header, \texttt{mcp\_servers}, \texttt{container}, \texttt{stream} \\
\sel{outputFormat:} \sel{effort:}
   & \cls{ClaudeOutputConfig}
   & \texttt{output\_config.format}, \texttt{output\_config.effort} \\
\bottomrule
\end{tabular}
\end{center}

Two fields do not live on the request even though new readers
often expect them there.  The \texttt{strict} flag is per-tool
and lives on \texttt{ClaudeTool>>strict:} (or on any
\texttt{ClaudeServerTool} subclass), because strictness is a
property of the tool schema the model is asked to honor, not of
the request that carries it.  A server-assigned
\texttt{requestId} is exposed response-side on
\texttt{ClaudeRequestContext>>requestId} and reached via
\texttt{message requestId}; the server fills it in, so it cannot
exist on a request that has not yet been sent.

Output encoding is a Pharo-specific convenience built on top of
the system prompt rather than a wire field.
\cls{ClaudeTextOutputFormat} combines two axes --- \texttt{encoding}
(\sel{ascii}, \sel{latin1}, \sel{utf8}) and \texttt{markup}
(\sel{plain}, \sel{microdown}, \sel{markdown}, \sel{json},
\sel{html}, \sel{latex}) --- into a prompt suffix that asks the
model to produce text in that shape.  The wire request carries
no \texttt{output\_format} field for text output; only
\texttt{json\_schema} has a dedicated wire surface.  Configure
the format on the request and the SDK appends the suffix during
\texttt{asJson}:

\begin{lstlisting}[language={}]
"Ask for UTF-8-encoded Markdown output."
request outputFormat: (ClaudeTextOutputFormat new
    encoding: #utf8;
    markup: #markdown;
    yourself).
\end{lstlisting}

Tool choice is a Dictionary that tells the model whether it may
use tools, must use some tool, or must use one specific named
tool.  The \texttt{auto} shape lets the model decide; the
\texttt{any} shape forces a tool call without pinning which; the
\texttt{tool} shape names the tool the model must call:

\begin{lstlisting}[language={}]
"Let the model decide (default)."
request toolChoice: (Dictionary new at: 'type' put: 'auto'; yourself).

"Force tool use without pinning which tool."
request toolChoice: (Dictionary new at: 'type' put: 'any'; yourself).

"Force one named tool."
request toolChoice: (Dictionary new
    at: 'type' put: 'tool';
    at: 'name' put: 'get_weather';
    yourself).
\end{lstlisting}

\cls{ClaudeThinkingParams} wraps the extended-thinking
configuration.  Four class-side factories cover the modes:
\sel{enabled:} turns thinking on with a token budget;
\sel{adaptive} lets the model decide; \sel{adaptiveWithDisplay:}
adds a display mode (\texttt{'summarized'} or \texttt{'full'});
\sel{disabled} turns it off.  Pass the resulting Dictionary to
\sel{thinking:} on the request:

\begin{lstlisting}[language={}]
"Always think, with a 10000-token budget."
request thinking: (ClaudeThinkingParams enabled: 10000) asDictionary.

"Let the model decide whether to think."
request thinking: ClaudeThinkingParams adaptive asDictionary.

"Adaptive with a display-mode hint."
request thinking: (ClaudeThinkingParams adaptiveWithDisplay: 'summarized') asDictionary.

"Thinking off."
request thinking: ClaudeThinkingParams disabled asDictionary.
\end{lstlisting}

\begin{sloppypar}
When thinking is enabled or adaptive, responses may contain
\cls{ClaudeThinkingBlock} and \cls{ClaudeRedactedThinkingBlock}
content blocks.  Redacted blocks appear when the API omits
thinking content for safety reasons; they carry opaque
\texttt{data} that callers must preserve when round-tripping
conversations so the next request retains the server's signed
prefix.  The full thinking surface is described in
\S\ref{sec:thinking}.
\end{sloppypar}

\begin{sloppypar}
Conversation turns are built through five builder methods.
\sel{addUserMessage:} appends a plain-text user turn;
\sel{addAssistantMessage:} appends a plain-text assistant turn
(useful for conversation replay); \sel{addAssistantResponse:}
appends a whole prior \cls{ClaudeMessage}, preserving its
content-block shape (thinking, citations, tool use all survive);
\sel{addToolResult:content:} and
\sel{addToolResult:content:isError:} append a user turn
containing a tool result paired to the matching tool-use id.
\end{sloppypar}

\subsection{Response Shape (ClaudeMessage)}

\cls{ClaudeMessage} is the typed response every successful
Messages call returns.  It holds seven instance variables ---
\texttt{id}, \texttt{type}, \texttt{turn}, \texttt{usage},
\texttt{stopDetails}, \texttt{executionContainer}, and
\texttt{requestContext} --- and exposes the familiar API shape
through forwarding accessors that delegate to the composite
value objects.  The public-accessor surface:

\begin{center}
\small
\begin{tabular}{@{}llp{6cm}@{}}
\toprule
\textbf{Accessor} & \textbf{Returns} & \textbf{Description} \\
\midrule
\sel{id}           & String         & Unique message identifier (instVar) \\
\sel{type}         & String         & Always \texttt{'message'} (instVar) \\
\sel{role}         & Symbol         & \texttt{\#assistant} for responses; delegated to \texttt{turn} \\
\sel{content}      & OrderedCollection of \cls{ClaudeContentBlock}
                   & Text, tool use, thinking blocks; delegated to \texttt{turn} \\
\sel{model}        & String         & Model ID used; delegated to \texttt{turn} \\
\sel{stopReason}   & Symbol         & See stop reason values below; delegated to \texttt{stopDetails} \\
\sel{stopSequence} & String or nil  & Matched stop sequence; delegated to \texttt{stopDetails} \\
\sel{usage}        & \cls{ClaudeUsage} & Token counts (instVar) \\
\bottomrule
\end{tabular}
\end{center}

The \texttt{stopReason} value is a Symbol the SDK normalizes
from the wire string: \texttt{\#end\_turn}, \texttt{\#tool\_use},
\texttt{\#max\_tokens}, \texttt{\#stop\_sequence},
\texttt{\#pause\_turn}, and \texttt{\#refusal}.  See
\S\ref{sec:stop-reasons} for per-value handling patterns.
Alongside the accessors, four key methods cover the
most-used response operations: \sel{asParam} converts the
message back to request-param shape for multi-turn replay;
\sel{accumulate:} builds the message incrementally from
streaming events (see \S\ref{sec:streaming}); \sel{textContent}
concatenates every \cls{ClaudeTextBlock} in \texttt{content}
into a single String; and \sel{toolUseBlocks} selects the
\cls{ClaudeToolUseBlock} entries for tool-dispatch code to
iterate.

The four composite value objects keep the response class
narrow.  \cls{ClaudeTurn} groups the role-and-content pair,
\cls{ClaudeStopDetails} holds the stop-reason breakdown,
\cls{ClaudeRequestContext} carries request-level metadata, and
\cls{ClaudeExecutionContainer} tracks the code-execution
container state when a request used skills or code execution:

\begin{center}
\small
\begin{tabular}{@{}llp{5.5cm}@{}}
\toprule
\textbf{Class} & \textbf{Responsibility} & \textbf{Key Fields} \\
\midrule
\cls{ClaudeTurn}             & Role and content grouping   & \texttt{role}, \texttt{content}, \texttt{model} \\
\cls{ClaudeStopDetails}      & Stop reason breakdown       & \texttt{stopReason}, \texttt{stopSequence}, \texttt{category}, \texttt{explanation} \\
\cls{ClaudeRequestContext}   & Request-level metadata      & \texttt{requestId}, \texttt{outputFormat} \\
\cls{ClaudeExecutionContainer} & Container execution context & \texttt{containerId}, \texttt{expiresAt}, \texttt{skills} \\
\bottomrule
\end{tabular}
\end{center}

\cls{ClaudeUsage} is the fifth composite and tracks the token
counts every response carries --- it is what cost-analysis and
cache-effectiveness code reads after a call:

\begin{center}
\small
\begin{tabular}{@{}llp{5.5cm}@{}}
\toprule
\textbf{Variable} & \textbf{Type} & \textbf{Description} \\
\midrule
\texttt{inputTokens}              & Integer          & Tokens in request \\
\texttt{outputTokens}             & Integer          & Tokens in response \\
\texttt{cacheCreationInputTokens} & Integer          & Prompt cache creation \\
\texttt{cacheReadInputTokens}     & Integer          & Prompt cache hits \\
\texttt{inferenceGeo}             & String or nil    & Inference region hint echoed by the server \\
\bottomrule
\end{tabular}
\end{center}

\subsection{Content Block Types}
\label{sec:content-blocks}

A response's \texttt{content} is a collection of
\cls{ClaudeContentBlock} instances, not a single string ---
Claude may return interleaved text, tool use, thinking, and
server-tool-result blocks in one message.  Each block is a
subclass with its own instance variables and its own
\sel{asJson} / \sel{fromJson:} pair.
\texttt{ClaudeContentBlock class>>fromJson:} dispatches on the
discriminator \texttt{type} field via a registry Dictionary
mapping type strings to subclasses; unknown types fall back to
a raw-JSON instance of the abstract class so forward-compatible
payloads never lose data.  Feature sections cross-reference this
taxonomy: Vision~(\S\ref{sec:vision}), PDF~(\S\ref{sec:pdf}),
Citations~(\S\ref{sec:citations}), Search
Results~(\S\ref{sec:search-results}),
Thinking~(\S\ref{sec:thinking}), and Tool
Use~(\S\ref{sec:tool-use}).  The primary hierarchy --- blocks
callers can both send and receive --- is:

\begin{center}
\small
\begin{tabular}{@{}llp{6.2cm}@{}}
\toprule
\textbf{Class} & \textbf{Wire Type} & \textbf{Key Fields} \\
\midrule
\cls{ClaudeContentBlock}            & (abstract)        & \texttt{type}, \texttt{rawJson} \\
\cls{ClaudeTextBlock}               & \texttt{text}     & \texttt{text}, \texttt{citations}, \texttt{cacheControl} \\
\cls{ClaudeToolUseBlock}            & \texttt{tool\_use} & \texttt{id}, \texttt{name}, \texttt{input} \\
\cls{ClaudeToolResultBlock}         & \texttt{tool\_result} & \texttt{toolUseId}, \texttt{content}, \texttt{isError}, \texttt{cacheControl} \\
\cls{ClaudeThinkingBlock}           & \texttt{thinking} & \texttt{thinking}, \texttt{signature} \\
\cls{ClaudeRedactedThinkingBlock}   & \texttt{redacted\_thinking} & \texttt{data} \\
\cls{ClaudeImageBlock}              & \texttt{image}    & \texttt{source} \\
\cls{ClaudeDocumentBlock}           & \texttt{document} & \texttt{source} (via \cls{ClaudeDocumentSource} hierarchy), \texttt{title}, \texttt{citations} \\
\cls{ClaudeContainerUploadBlock}    & \texttt{container\_upload} & \texttt{fileId} (derived from \texttt{rawJson}) \\
\cls{ClaudeCompactionBlock}         & \texttt{compaction} & \texttt{content} (server-emitted summary) \\
\cls{ClaudeMCPToolUseBlock}         & \texttt{mcp\_tool\_use} & \texttt{id}, \texttt{name}, \texttt{serverName}, \texttt{input} \\
\cls{ClaudeMCPToolResultBlock}      & \texttt{mcp\_tool\_result} & \texttt{toolUseId}, \texttt{content}, \texttt{isError} \\
\bottomrule
\end{tabular}
\end{center}

\cls{ClaudeDocumentBlock} sends documents to the API.  The wire
format uses type \texttt{"document"} with a \texttt{source}
object whose structure varies by source type; the
\cls{ClaudeDocumentSource} hierarchy covers five shapes:

\begin{center}
\small
\begin{tabular}{@{}llp{5cm}@{}}
\toprule
\textbf{Class} & \textbf{Source Type} & \textbf{Key Fields} \\
\midrule
\cls{ClaudeBase64DocumentSource}  & \texttt{base64}  & \texttt{data}, \texttt{mediaType} \\
\cls{ClaudeTextDocumentSource}    & \texttt{text}    & \texttt{data} \\
\cls{ClaudeContentDocumentSource} & \texttt{content} & \texttt{content} (array of blocks) \\
\cls{ClaudeURLDocumentSource}     & \texttt{url}     & \texttt{url} \\
\cls{ClaudeFileDocumentSource}    & \texttt{file}    & \texttt{fileId} \\
\bottomrule
\end{tabular}
\end{center}

\begin{sloppypar}
Four class-side factories on \cls{ClaudeDocumentBlock} ---
\sel{fromBase64:}, \sel{fromFileId:}, \sel{fromText:},
\sel{fromURL:} --- pick the right source subclass so callers
rarely touch \cls{ClaudeDocumentSource} directly; on the
response side, \texttt{ClaudeDocumentSource class>>fromJson:}
dispatches on the \texttt{type} field and returns the matching
concrete source.  \cls{ClaudeTextBlock}'s \texttt{citations}
collection uses the \cls{ClaudeCitation} shape described in
\S\ref{sec:citations}.
\end{sloppypar}

Five server-managed blocks appear only in responses and are
read-only from a client perspective.  The SDK deserializes them
so conversations round-trip, but application code does not
build them:

\begin{center}
\small
\begin{tabular}{@{}lp{9cm}@{}}
\toprule
\textbf{Class} & \textbf{Description} \\
\midrule
\cls{ClaudeServerToolUseBlock}        & Server-managed tool invocation (wire type \texttt{'server\_tool\_use'}) \\
\cls{ClaudeWebSearchResultBlock}      & Web search result (wire type \texttt{'web\_search\_tool\_result'}) \\
\cls{ClaudeWebFetchResultBlock}       & Web fetch result (wire type \texttt{'web\_fetch\_tool\_result'}) \\
\cls{ClaudeCodeExecutionResultBlock}  & Code execution result (wire type \texttt{'code\_execution\_tool\_result'}) \\
\cls{ClaudeBashResultBlock}           & Server Bash result (wire type \texttt{'bash\_code\_execution\_tool\_result'}); fields \texttt{toolUseId}, \texttt{returnCode}, \texttt{stdout}, \texttt{stderr} \\
\cls{ClaudeTextEditorResultBlock}     & Server text-editor result (wire type \texttt{'text\_editor\_code\_execution\_tool\_result'}); fields \texttt{toolUseId}, \texttt{rawJson} (content sub-dict shape varies by editor command: \texttt{view} / \texttt{create} / \texttt{str\_replace}) \\
\cls{ClaudeToolSearchResultBlock}     & Tool-discovery result (wire type \texttt{'tool\_search\_tool\_result'}); fields \texttt{toolUseId}, \texttt{toolReferences} \\
\bottomrule
\end{tabular}
\end{center}

Each server-result block stores its raw JSON payload so fields
the server adds in the future survive the round trip.

\subsection{Per-Request Options}

The durable request controls --- metadata, cache control,
service tier, beta flags, MCP server list, container, and the
streaming mode bit --- live on \cls{ClaudeRequestOptions} and
are set through the forwarded setters on
\cls{ClaudeMessageRequest} (see the Request Parameters tables
above).  Those options travel with the request on the wire; the
SDK never treats them as ad-hoc closures.

A per-call override that doesn't belong on the wire --- a
timeout bump for one particular long-running request, a debug
log block scoped to one call --- is a property of the client,
not of the request.  The preferred idiom is a dedicated
\cls{ClaudeClient} configured with the override, used for the
one call and then discarded.  Clients are cheap and stateless
apart from their retry counter, so spinning up a second client
with a different timeout is the safe way to scope the change
--- a shared client never ends up stuck in the overridden state
if the call raises:

\begin{lstlisting}[language={}]
"Dedicated client with a tighter timeout for one slow call."
tightClient := ClaudeClient
    apiKey: client apiKey
    baseUrl: client baseUrl
    timeout: 30.
response := tightClient sendMessage: request.
\end{lstlisting}

The same pattern is the right shape for a parallel fan-out that
wants a per-arm timeout: build one \cls{ClaudeClient} per arm
with the timeout baked in, let each arm own its client, and the
shared default client is never mutated.

When building a second client is genuinely inconvenient ---
deep inside a helper that receives the client as an argument,
say --- mutate the setter inside an \sel{ensure:} block so the
original value is restored even if the call raises:

\begin{lstlisting}[language={}]
"Ad-hoc override on the shared client, restored on every exit path."
| originalTimeout |
originalTimeout := client timeout.
client timeout: 30.
[ response := client sendMessage: request ]
    ensure: [ client timeout: originalTimeout ].
\end{lstlisting}

For the canonical first-contact example see
\S\ref{sec:overview}; for the tokens-only counterpart see
\S\ref{sec:count-tokens}; for the exception hierarchy every
call can raise see \S\ref{sec:errors}; for extended thinking
and tool use that ride on a Messages request see
\S\ref{sec:thinking} and \S\ref{sec:tool-use}; for streaming
see \S\ref{sec:streaming}.

\medskip
\noindent{\small\textsf{Status: \emph{Available}.\quad
Anthropic docs: \url{https://docs.claude.com/en/api/messages}.\quad
Python SDK peer:
\url{https://docs.anthropic.com/en/api/messages}.}}

\section{Messages Examples}
\label{sec:messages-examples}

\begin{sloppypar}
This section is the pattern catalog that accompanies
\S\ref{sec:messages} Messages API.  Where \S\ref{sec:messages}
describes the wire surface and its Pharo bindings,
\S\ref{sec:messages-examples} collects the runnable
conversation shapes a reader may reach for: multi-turn
conversations, conversation replay, and tool-result
round-tripping.  Every pattern is an ordinary
\cls{ClaudeMessageRequest} built through the five builder
selectors (\sel{addUserMessage:}, \sel{addAssistantMessage:},
\sel{addAssistantResponse:}, \sel{addToolResult:content:},
\sel{addToolResult:content:isError:}); the SDK's job is to
serialize that request and return a typed
\cls{ClaudeMessage}.
\end{sloppypar}

The runnable below is a multi-turn conversation against
\cls{ClaudeMockServer}, which makes the pattern testable in the
eval server without a live API key or network.  The first turn
is a user question; the mock returns an assistant answer; the
example appends that full response to the request via
\sel{addAssistantResponse:}, adds a follow-up user turn, and
sends the request a second time.  The accumulated messages
collection after both round-trips contains four turns (user,
assistant, user, assistant), which is the predicate the example
evaluates to:

\begin{lstlisting}[language={}]
| server client request response |
"Mock the same canned response for both POSTs; the example only checks turn bookkeeping."
server := ClaudeMockServer startOn: 0.
[ server
    respondTo: '/v1/messages'
    with: '{"id":"msg_01","type":"message","role":"assistant","content":[{"type":"text","text":"2 + 2 is 4."}],"model":"claude-sonnet-4-5","stop_reason":"end_turn","stop_sequence":null,"usage":{"input_tokens":12,"output_tokens":7,"cache_creation_input_tokens":0,"cache_read_input_tokens":0}}'.
  client := ClaudeClient
    apiKey: 'test-key'
    baseUrl: 'http://localhost:' , server port printString
    timeout: 5.
  "Turn 1: build the request and send the first user question."
  request := ClaudeMessageRequest new
    model: ClaudeModel sonnet45;
    maxTokens: 256;
    addUserMessage: 'What is 2 + 2?';
    yourself.
  response := client sendMessage: request.
  "Append the full assistant response; addAssistantResponse: preserves every content block."
  request addAssistantResponse: response.
  "Turn 2: add a follow-up user question and reuse the client."
  request addUserMessage: 'And what is that times 3?'.
  response := client sendMessage: request.
  "The request now has four turns: user, assistant, user, assistant."
  ^ request messages size ]
    ensure: [ server stop ]
"=> 4"
\end{lstlisting}

\sel{addAssistantResponse:} takes a whole \cls{ClaudeMessage}
and appends it as an assistant turn, calling \sel{asParam} on
the message to produce the request-side shape.  Each
\cls{ClaudeContentBlock} subclass answers its own \sel{asParam},
so text, thinking, citations, and redacted-thinking blocks all
survive the round trip correctly --- preserving redacted
thinking is critical when the server's signed prefix would
otherwise be lost.  \sel{addAssistantMessage:} is the
plain-String variant: when replaying a conversation from a
transcript, the reader often has the assistant text as a String
and not as a \cls{ClaudeMessage}, and this selector wraps it
back into a turn.  Tool-use loops use
\sel{addToolResult:content:} (or the \texttt{isError:} variant
to mark a failed tool call) to pair a result with its
tool-use id.

Further patterns --- prefill with an assistant turn, vision
with an image content block, JSON-schema output, and full
tool-use loops --- live with the feature sections they
exercise: \S\ref{sec:thinking} for extended thinking,
\S\ref{sec:vision} for vision, \S\ref{sec:citations} for
citations, \S\ref{sec:tool-use} for client-side tool use,
\S\ref{sec:server-tools} for server-executed tools,
\S\ref{sec:structured-outputs} for JSON schema, and
\S\ref{sec:errors} for the error-handling shapes that wrap
every call.

\medskip
\noindent{\small\textsf{Status: \emph{Available}.\quad
Anthropic docs: \url{https://docs.claude.com/en/api/messages-examples}.}}

\section{Count Tokens}
\label{sec:count-tokens}

Counting tokens before you send a message lets the caller quote
cost up front and gate on context-window headroom without paying
for a round-trip that would have exceeded the budget.  The
count\_tokens endpoint accepts the same request shape as
\texttt{sendMessage:} --- the same model, messages, system prompt,
and tools --- so the count you receive is the same count the
billing path would meter a moment later.  Applications use this
for pre-flight pricing (``this conversation will cost N tokens''),
UX warnings when a long turn is approaching the model's input
cap, and batch-job gating where you would rather skip a request
than overrun a quota.

\begin{sloppypar}
\texttt{ClaudeClient>>countTokens:} sends \texttt{POST
/v1/messages/count\_tokens} and returns a
\texttt{ClaudeMessageTokensCount}, a value object carrying one
field: \texttt{inputTokens} (Integer) --- the token count for the
given request parameters.  The response object has no output-token
or cached-token breakdown; that information only exists once the
model has actually generated a response.
\end{sloppypar}

The example below builds a request the same way
\texttt{sendMessage:} takes one, hands it to \texttt{countTokens:},
and reads the count off the result.  It is an end-to-end pre-flight
check: if \texttt{count inputTokens} exceeds the caller's budget,
skip the send entirely; otherwise hand the same \texttt{params}
object to \texttt{sendMessage:}.

\begin{lstlisting}[language={},title={Count tokens}]
| params count |
"Build a request the same way sendMessage: would."
params := ClaudeMessageRequest new
    model: ClaudeModel opus47;
    maxTokens: 1024;
    addUserMessage: 'How many tokens is this?';
    yourself.

"Send the count-tokens request; surface any API error directly."
[ count := client countTokens: params ]
  on: ClaudeApiError do: [ :e | ^ e signal ].

"count is a ClaudeMessageTokensCount; read the single field."
count inputTokens  "=> 14"
\end{lstlisting}

For the sibling endpoint that actually consumes the counted tokens
see \S\ref{sec:messages}.

\medskip
\noindent{\small\textsf{Status: \emph{Available}.\quad
Anthropic docs: \url{https://docs.claude.com/en/api/messages-count-tokens}.\quad
Python SDK peer:
\url{https://github.com/anthropics/anthropic-sdk-python\#token-counting}.}}

\section{Models API --- List Models}
\label{sec:models-list}

The Models API tells an application what Claude variants the
account can actually call, with their capability limits --- useful
for model discovery in a UI, for capability gating (``skip this
model if \texttt{maxInputTokens} is below the prompt I need to
send''), and for spotting new aliases when Anthropic publishes a
refreshed family.  \texttt{ClaudeClient>>listModels} sends
\texttt{GET /v1/models} and returns an \texttt{OrderedCollection}
of \texttt{ClaudeModelInfo}; each entry carries the model ID,
display name, token limits, and creation timestamp.

The example below lists every available model and prints its
display name alongside its ID --- the shape of a discovery UI or a
capability audit.  \texttt{listModels} is synchronous; a
paginated variant, \texttt{listModels:after:}, returns a
\texttt{Dictionary} with \texttt{data}, \texttt{has\_more},
\texttt{first\_id}, and \texttt{last\_id} when cursor control
matters.

\begin{lstlisting}[language={},title={List available models}]
| models |
"Fetch every model the account can call; surface API errors directly."
[ models := client listModels ]
  on: ClaudeApiError do: [ :e | ^ e signal ].

"Walk the collection; each entry is a ClaudeModelInfo."
models do: [ :info |
    Transcript show: info displayName, ' (', info id, ')'; cr ].
\end{lstlisting}

For code that wants to name a model at a call site without
hard-coding a string, \texttt{ClaudeModel} provides class-side
factory methods returning the model ID strings.  Applications may
also pass any string the API accepts; the constants are
conveniences that document which model family a call targets.  The
factories and the strings they return:

\begin{lstlisting}[language={},title={Model constants}]
ClaudeModel haiku35     "=> 'claude-3-haiku-20240307'"
ClaudeModel haiku45     "=> 'claude-haiku-4-5-20251001'"
ClaudeModel sonnet45    "=> 'claude-sonnet-4-5-20250929'"
ClaudeModel sonnet45v   "=> 'claude-sonnet-4-5-20250929'"
ClaudeModel sonnet46    "=> 'claude-sonnet-4-6'"
ClaudeModel opus4       "=> 'claude-opus-4-20250514'"
ClaudeModel opus4v      "=> 'claude-opus-4-20250514'"
ClaudeModel opus46      "=> 'claude-opus-4-6'"
ClaudeModel opus46v     "=> 'claude-opus-4-6'"
ClaudeModel opus47      "=> 'claude-opus-4-7'"
ClaudeModel opus47v     "=> 'claude-opus-4-7'"
\end{lstlisting}

The \texttt{v} suffix historically distinguished dated model IDs
from floating aliases.  In the current table, non-\texttt{v} and
\texttt{v} selectors return the same string for the newer models
(\texttt{opus46}, \texttt{opus47}, \texttt{sonnet46}); the suffix
exists for future-proofing when Anthropic publishes a dated
variant.  For older models (\texttt{opus4}, \texttt{sonnet45}),
both forms already resolve to a dated string --- there is no
functional floating alias.

For the single-model sibling endpoint see \S\ref{sec:models-get};
for the scope of the Models API surface in this SDK see
\S\ref{sec:scope}.

\medskip
\noindent{\small\textsf{Status: \emph{Available}.\quad
Anthropic docs: \url{https://docs.claude.com/en/api/models-list}.\quad
Python SDK peer: \url{https://github.com/anthropics/anthropic-sdk-python\#list-models}.}}

\section{Models API --- Get a Model}
\label{sec:models-get}

When an application already knows a model ID and needs only that
model's capability details --- maximum input tokens, maximum output
tokens, the human-readable display name --- fetching the single
record is cheaper than listing the catalog and filtering.
\texttt{ClaudeClient>>getModel:} sends \texttt{GET
/v1/models/\{id\}} and returns a single \texttt{ClaudeModelInfo};
an unknown ID raises \texttt{ClaudeNotFoundError} (404), which a
capability-gating path typically treats as ``model unavailable''
rather than a hard failure.

The example below retrieves one model by ID, handles the
unknown-ID case with the typed 404, and then reads the
capability fields the application cares about.  The
\texttt{on: ClaudeNotFoundError} idiom is the right shape for a
capability check: other errors (network, 5xx, auth) propagate.

\begin{lstlisting}[language={},title={Get specific model info}]
| info |
"Retrieve capability details for a specific model;
 catch 404 as 'model unavailable' and let other errors propagate."
[ info := client getModel: 'claude-sonnet-4-5' ]
  on: ClaudeNotFoundError do: [ :e |
    Transcript show: 'Unknown model'; cr. ^ nil ].

"Read the capability fields off the ClaudeModelInfo response."
info maxInputTokens  "=> 200000"
info maxTokens       "=> 8192"
\end{lstlisting}

\texttt{ClaudeModelInfo} is the response type for both
\texttt{listModels} and \texttt{getModel:}.  It carries six fields:

\begin{center}
\begin{tabular}{@{}llp{5.5cm}@{}}
\toprule
\textbf{Variable} & \textbf{Type} & \textbf{Description} \\
\midrule
\texttt{id}              & String   & Model identifier \\
\texttt{displayName}     & String   & Human-readable name \\
\texttt{maxInputTokens}  & Integer  & Maximum input context length \\
\texttt{maxTokens}       & Integer  & Maximum output tokens \\
\texttt{type}            & String   & Always \texttt{'model'} \\
\texttt{createdAt}       & String   & ISO-8601 creation timestamp \\
\bottomrule
\end{tabular}
\end{center}

For the list endpoint that returns an
\texttt{OrderedCollection} of these records see
\S\ref{sec:models-list}.

\medskip
\noindent{\small\textsf{Status: \emph{Available}.\quad
Anthropic docs: \url{https://docs.claude.com/en/api/models}.\quad
Python SDK peer: \url{https://github.com/anthropics/anthropic-sdk-python\#retrieve-model}.}}

\section{Handling Stop Reasons}
\label{sec:stop-reasons}

Every Messages response carries a \texttt{stopReason} --- a Symbol
that tells the caller why Claude stopped generating.  Applications
dispatch on it to continue a tool-use loop, resume a paused turn,
surface a natural completion to the user, or fall through to
telemetry when the model ran into a limit.  The stop reason is the
first thing an agent reads off a response; it is how the agent
decides what to do next.

\begin{sloppypar}
The wire protocol actually uses two value domains and the SDK
preserves them.  The top-level \texttt{stop\_reason} field from the
response is parsed by \texttt{ClaudeStopDetails} class>>\texttt{fromJson:}
and surfaced as \texttt{ClaudeMessage>>stopReason}, a Symbol.  The
separate \texttt{stop\_details.type} sub-object (currently used for
the \texttt{'refusal'} category with its own \texttt{category} and
\texttt{explanation} fields) is accessed through
\texttt{ClaudeMessage>>stopDetails} and the \texttt{category} /
\texttt{explanation} accessors --- it does \emph{not} appear as a
\texttt{stopReason} value.  The table below covers the five
Symbol values the \texttt{stop\_reason} field actually emits:
\end{sloppypar}

\begin{center}
\begin{tabular}{@{}lp{9.5cm}@{}}
\toprule
\textbf{stopReason} & \textbf{Meaning} \\
\midrule
\texttt{\#end\_turn}     & Natural completion \\
\texttt{\#max\_tokens}   & Hit \texttt{maxTokens} cap or model max \\
\texttt{\#stop\_sequence} & Hit a \texttt{stopSequences} entry \\
\texttt{\#tool\_use}     & Model invoked one or more tools \\
\texttt{\#pause\_turn}   & Long-running turn paused; resume with the same messages \\
\bottomrule
\end{tabular}
\end{center}

Refusals arrive as \texttt{stopReason = \#end\_turn} plus a
populated \texttt{stopDetails} sub-object whose \texttt{category}
names the safety category (for example,
\texttt{'violence'}) and whose \texttt{explanation} carries a
short rationale.  Code that needs to distinguish a refusal from an
ordinary natural completion reads \texttt{response stopDetails
category} after checking \texttt{stopReason = \#end\_turn} ---
both paths use the same \texttt{\#end\_turn} stop reason on the
wire.

Applications that need to branch on \texttt{stopReason} match on
the Symbol directly --- there is no wrapper class or visitor
pattern.  A typical first-contact dispatch covers the three
branches an agent loop actually cares about: complete response,
tool-use continuation, and paused-turn resumption.  The runnable
below starts with a nil-guard: \texttt{ClaudeMessage>>stopReason}
returns nil when \texttt{stopDetails} has not been populated (an
in-flight streaming response, a manually constructed message),
and silent nil-dispatch is not a shape any caller wants:

\begin{lstlisting}[language={},title={Dispatch on stopReason}]
| response reason |
response := client sendMessage: request.
reason := response stopReason.

"Nil-guard: an incomplete response has no stopReason yet."
reason ifNil: [ ^ nil ].

"Natural completion: hand the text back to the caller."
reason = #end_turn ifTrue: [ ^ response textContent ].

"Tool use: the model invoked tools; run them and loop."
reason = #tool_use ifTrue: [ ^ self handleToolUse: response ].

"Paused turn: long-running work suspended; resume with same messages."
reason = #pause_turn ifTrue: [ ^ self resumeWithSameMessages ].

"Fall-through: #max_tokens or #stop_sequence --- telemetry / surface."
^ self recordUnexpectedStop: reason
\end{lstlisting}

For the refusal case --- and for other non-2xx failures that
never produce a response in the first place --- see
\S\ref{sec:errors}; for the \texttt{\#tool\_use} branch and how a
tool-use loop continues see \S\ref{sec:tool-use}; for the
\texttt{\#pause\_turn} branch and the Extended Thinking flow that
drives it see \S\ref{sec:thinking}.

\medskip
\noindent{\small\textsf{Status: \emph{Available}.\quad
Anthropic docs: \url{https://docs.claude.com/en/api/handling-stop-reasons}.}}

\section{Errors}
\label{sec:errors}

The Messaging SDK turns every non-2xx HTTP response from the
Messages API into a typed Smalltalk exception.  Rather than carrying
a status-code field on a single error class, the SDK dispatches at
the response boundary: a 401 becomes a \texttt{ClaudeAuthError},
a 429 becomes a \texttt{ClaudeRateLimitError}, a 529 becomes a
\texttt{ClaudeOverloadedError}.  Callers write
\texttt{[ ... ] on: ClaudeAuthError do: [:e | ...]} and let
everything else propagate --- the idiomatic Smalltalk analog of the
Python SDK's \texttt{APIStatusError} hierarchy.

\texttt{ClaudeApiError} is the abstract base, a subclass of
\texttt{Error}.  Ten concrete subclasses cover the status codes
the API actually returns:

\begin{center}
\begin{tabular}{@{}l l c p{5.6cm}@{}}
\toprule
\textbf{Class} & \textbf{HTTP Status} & \textbf{Retryable} & \textbf{When} \\
\midrule
\texttt{ClaudeApiError}            & (abstract)       & ---  & Base class for all API errors \\
\texttt{ClaudeInvalidRequestError} & 400              & no   & Malformed or invalid request \\
\texttt{ClaudeAuthError}           & 401              & no   & Invalid API key \\
\texttt{ClaudeBillingError}        & 402              & no   & Billing issue \\
\texttt{ClaudePermissionError}     & 403              & no   & Insufficient permissions \\
\texttt{ClaudeNotFoundError}       & 404              & no   & Unknown model or endpoint \\
\texttt{ClaudeTimeoutError}        & 408 / 504        & yes  & Request or gateway timeout \\
\texttt{ClaudeRateLimitError}      & 429              & yes  & Rate limit exceeded; honors \texttt{Retry-After} \\
\texttt{ClaudeServerError}         & 500              & yes  & Internal server error \\
\texttt{ClaudeOverloadedError}     & 529              & yes  & API overloaded (Anthropic-specific) \\
\texttt{ClaudeNetworkError}        & 599 (synthetic)  & no   & Connection refused, TLS handshake, DNS --- a refused connection stays refused \\
\bottomrule
\end{tabular}
\end{center}

Every instance carries the status code, the Anthropic error message,
the request ID from the \texttt{request-id} response header, the raw
JSON body for logging, and an optional \texttt{retryAfter} parsed
from the \texttt{retry-after} header on 429 responses.  The
class-side factory \texttt{fromStatusCode:json:requestId:} is what
\texttt{ClaudeClient} calls internally; in application code you
almost never construct these directly --- you catch them.

The example below demonstrates the dispatch end-to-end.  It calls
the factory with a 401 status code, signals the resulting
exception, and shows a handler that catches the specific subclass
(\texttt{ClaudeAuthError}) the factory routed to.  The same shape
is what \texttt{ClaudeClient} produces on a real auth failure, so
the application code that handles it looks identical:

\begin{lstlisting}[language={}]
| caught |
caught := [
    "Build the error and raise it (the SDK does this internally for any non-2xx response)."
    (ClaudeApiError
        fromStatusCode: 401
        json: '{"error":{"message":"bad key"}}'
        requestId: 'req_1') signal ]
    on: ClaudeAuthError
    do: [ :e | e class name ].  "Handler catches the specific subclass."
caught
"=> #ClaudeAuthError"
\end{lstlisting}

Streaming responses get the same exception hierarchy.  The
Messages API uses Server-Sent Events for incremental token
delivery (see \S\ref{sec:streaming}); an in-flight stream can
deliver an SSE \texttt{error} event after the response has started
--- a server-side overload, a model timeout, a 529 spike.  The SDK
routes those events through \texttt{fromSseErrorData:}, populating
the synthetic status code from the Anthropic \texttt{error.type}
string so \texttt{isRetryable} and \texttt{printOn:} behave
identically whether the error arrived on the status line or in a
streaming payload.  One \texttt{on: ClaudeApiError do: [...]}
handler covers both paths.

The SDK retries the retryable failures (the \emph{Retryable}
column above) automatically --- five retries by default, with
exponential backoff from 0.5\,s to 8\,s and \textpm 25\% jitter.
When the response carries a \texttt{Retry-After-Ms} or
\texttt{Retry-After} header, the SDK uses the header value instead
of the computed delay, clamped to 30\,s to defend against
pathological values.  Tests that exercise the non-retry path must
explicitly set \texttt{maxRetries: 0} or they will wait out the
full backoff sequence; calls from the UI thread must be wrapped in
a forked Process because the backoff delay blocks the current
process.

Two divergences from the Python SDK are worth knowing.  HTTP 529
is Anthropic-specific and not IANA-assigned --- the SDK treats it
as a retryable 5xx.  HTTP 599 is synthetic; the Python SDK raises
\texttt{APIConnectionError} without a status code, while the Pharo
SDK keeps everything inside the \texttt{ClaudeApiError} hierarchy
so a single \texttt{on: ClaudeApiError do: [...]} catches both
wire-protocol errors and connection errors uniformly.

For the rate-limit monitoring headers surfaced on
\texttt{ClaudeRateLimitError} see \S\ref{sec:rate-limits}; for
how mid-stream SSE errors are routed see \S\ref{sec:streaming}.

\medskip
\noindent{\small\textsf{Status: \emph{Available}.\quad
Anthropic docs: \url{https://docs.claude.com/en/api/errors}.\quad
Python SDK peer:
\url{https://github.com/anthropics/anthropic-sdk-python\#handling-errors}.}}

\section{API Versions}
\label{sec:api-versions}

Every request to the Messages API carries an \texttt{anthropic-version}
header that pins the wire format.  Anthropic guarantees that requests
sent with a given version keep the same request and response shape
indefinitely, so a client that pins a version keeps working across
Anthropic's platform changes until it opts into a newer one.  The
baseline is \texttt{2023-06-01}, which is what every first-party SDK
ships with; new versions land when Anthropic changes response shapes
or introduces incompatible defaults.

In Pharo the version lives on \texttt{ClaudeClient}'s
\texttt{apiVersion} slot, initialized to \texttt{'2023-06-01'} on
construction and sent as the \texttt{anthropic-version} header by
every HTTP-client factory (\texttt{newHttpClient},
\texttt{newFileHttpClient}, \texttt{newSkillHttpClient}).  Override it
with \texttt{apiVersion:} on a per-client basis when you need to opt
into a new wire version during migration.

This section is reference material; no runnable example applies.

\medskip
\noindent{\small\textsf{Status: \emph{Available}.\quad
Anthropic docs: \url{https://docs.claude.com/en/api/versioning}.}}

\section{Beta Headers}
\label{sec:beta-headers}

Some Messages API features sit behind an \texttt{anthropic-beta}
header that gates access to not-yet-stable surface: new input shapes
(PDFs, extended-output models), new endpoints (Files API, Skills API,
MCP client), or experimental response semantics.  Each beta is a
string identifier like \texttt{pdfs-2024-09-25} or
\texttt{files-api-2025-04-14}; callers can opt into several at once
by sending a comma-joined header, and Anthropic lists the active
identifiers on its beta-headers documentation page.

Pharo handles beta headers two ways, depending on whether the beta
is an SDK implementation detail or an application-level choice.
\texttt{ClaudeClient} owns the automatic path: \texttt{newFileHttpClient}
sends the Files API beta on every Files request and
\texttt{newSkillHttpClient} sends the Skills beta on every Skills
request, so Files and Skills callers never touch beta headers
themselves.  For per-request betas --- the MCP client, PDF support,
extended-output models --- callers set \texttt{betas:} on
\texttt{ClaudeRequestOptions} (directly, or via
\texttt{ClaudeMessageRequest} which delegates).  At send time,
\texttt{newHttpClientFor:} reads the collection off the request
options, joins it with commas, and emits it as the
\texttt{anthropic-beta} header.

\begin{sloppypar}
The SDK ships a typed catalog of the known beta identifiers as
\texttt{ClaudeBetaHeader}, a class intended for class-side use.
Class-side messages return the wire String:
\texttt{ClaudeBetaHeader pdfs} answers
\texttt{'pdfs-2024-09-25'},
\texttt{ClaudeBetaHeader structuredOutputs} answers
\texttt{'structured-outputs-2025-11-13'}, and so on.  Typos are
surfaced when the send is evaluated (as an unknown selector /
\texttt{doesNotUnderstand:}) instead of landing only as a server
400, and \texttt{ClaudeBetaHeader allKnown} returns
the full list as an \texttt{OrderedCollection} for docs generation,
validation, and test parity.  The catalog is pure discovery:
\texttt{ClaudeRequestOptions>>betas:} still accepts an arbitrary
collection of Strings, so unknown-to-the-SDK identifiers remain
usable without waiting on a catalog update.
\end{sloppypar}

The example below demonstrates the opt-in path in both spellings.
It builds a request options object, sets two beta identifiers, and
asks it for the header value that \texttt{ClaudeClient} would emit
on the wire --- the same value that \texttt{newHttpClientFor:}
would place on the \texttt{anthropic-beta} header.  The first body
matches \texttt{ClaudeRequestOptions}'s class-comment doctest
byte-for-byte; the second shows the catalog-aware form and
produces the identical wire value:

\begin{lstlisting}[language={}]
"Build an options object with two active betas and read back the wire value."
ClaudeRequestOptions new
    betas: #('pdfs-2024-09-25' 'output-128k-2025-02-19');
    betaHeader
"=> 'pdfs-2024-09-25,output-128k-2025-02-19'"
\end{lstlisting}

\begin{lstlisting}[language={}]
"Same two betas via the typed catalog -- identical wire value."
ClaudeRequestOptions new
    betas: { ClaudeBetaHeader pdfs. ClaudeBetaHeader output128k };
    betaHeader
"=> 'pdfs-2024-09-25,output-128k-2025-02-19'"
\end{lstlisting}

To use this in an application, call \texttt{betas:} directly on the
\texttt{ClaudeMessageRequest} before \texttt{sendMessage:} --- the
request owns a \texttt{ClaudeRequestOptions} instance from
construction and the convenience setter writes through to it.  For
the less common cases that \texttt{betas:} does not cover, mutate
\texttt{request requestOptions} directly.  The Files and Skills APIs
require no such wiring --- their beta headers are part of the
endpoint's identity and the SDK adds them automatically via
\texttt{ClaudeBetaHeader filesApi} and
\texttt{ClaudeBetaHeader skills}.

\medskip
\noindent{\small\textsf{Status: \emph{Available}.\quad
Anthropic docs: \url{https://docs.claude.com/en/api/beta-headers}.}}

\section{Rate Limits}
\label{sec:rate-limits}

Anthropic enforces rate limits per API key on three axes: requests
per minute (RPM), input tokens per minute (ITPM), and output tokens
per minute (OTPM).  The limits scale with the account tier, and
every response carries monitoring headers
(\texttt{anthropic-ratelimit-requests-remaining},
\texttt{anthropic-ratelimit-tokens-remaining}, and companions) that
tell the caller how close they are to the next 429.  When a limit is
hit, the API returns a 429 with a \texttt{retry-after} hint; the
Pharo SDK surfaces this as \texttt{ClaudeRateLimitError}
(\S\ref{sec:errors}) and honors the hint automatically on retry.

For applications that make many concurrent requests on the same API
key, reactive 429 handling is not enough --- the retry-after delay
blocks the calling process, and several callers racing to retry can
stampede the budget right back to zero.  The SDK ships
\texttt{ClaudeRateLimiter}, a token-bucket limiter that coordinates
request rates across every \texttt{ClaudeClient} sharing a key.  The
bucket starts full (five tokens by default, one per permitted burst
request), drains one token per call to \texttt{acquire}, and refills
at \texttt{requestsPerMinute / 60} tokens per second.  When the
bucket is empty, callers block in \texttt{acquire} until a token
becomes available or the 30-second timeout fires.

\begin{sloppypar}
The example below demonstrates constructing a limiter at a chosen
rate.  The body matches \texttt{ClaudeRateLimiter}'s class-comment
doctest byte-for-byte --- the printed form is the limiter's own
\texttt{printOn:}, which reports its configured RPM and burst size:
\end{sloppypar}

\begin{lstlisting}[language={}]
"Build a limiter sized for 50 requests per minute and read its printed form."
ClaudeRateLimiter new requestsPerMinute: 50
"=> a ClaudeRateLimiter(50 rpm, 5 burst)"
\end{lstlisting}

To use the limiter with a client, attach it via \texttt{rateLimiter:}
on \texttt{ClaudeClient}; \texttt{sendMessage:} then calls
\texttt{acquire} before dispatching each HTTP request, blocking when
the bucket runs dry.  The class-side factories
\texttt{forKey:} and \texttt{forKey:model:} return a shared limiter
keyed on an API key (and optionally the model name), so every client
built against the same credentials automatically contends on the
same bucket --- the default model-aware factory picks 20 RPM for
Opus, 50 RPM for Haiku, and 40 RPM elsewhere.

\begin{sloppypar}
For reactive observation of the live limit state, every
\texttt{sendMessage:} response carries a typed
\texttt{ClaudeResponseMetadata} attached as
\texttt{ClaudeMessage>>responseMetadata} (with a
\texttt{rateLimits} convenience reader for shorter call chains).
The metadata holds eleven nil-tolerant slots: \texttt{requestId}
echoed from the \texttt{request-id} header, the three-axis
\texttt{anthropic-ratelimit-*} counters
(\texttt{rateLimitRequestsLimit},
\texttt{rateLimitRequestsRemaining},
\texttt{rateLimitRequestsReset},
\texttt{rateLimitInputTokensLimit},
\texttt{rateLimitInputTokensRemaining},
\texttt{rateLimitInputTokensReset},
\texttt{rateLimitOutputTokensLimit},
\texttt{rateLimitOutputTokensRemaining},
\texttt{rateLimitOutputTokensReset}), and the optional
\texttt{retryAfter} hint that 429 responses set per RFC 7231 §7.1.3.
The integer counters parse via \texttt{asInteger}; reset timestamps
stay as the raw ISO-8601 String the server sends so callers can
parse on demand without losing precision.  The
\texttt{hasRateLimits} predicate answers \texttt{true} when any of
the nine ratelimit slots is populated; \texttt{requestId} and
\texttt{retryAfter} never count toward it.  The same metadata
attaches to \texttt{ClaudeApiError} subclasses on the failure path
so \texttt{ClaudeRateLimitError} handlers can read live counters and
the reset timestamp to pick a smarter backoff than
\texttt{Retry-After} alone.
\end{sloppypar}

\begin{lstlisting}[language={}]
"Read the live ratelimit counters off a successful send."
| client request response |
client := ClaudeSDKExampleSupport resolveClient.
request := ClaudeMessageRequest new
    model: ClaudeModel sonnet46;
    maxTokens: 64;
    addUserMessage: 'ping';
    yourself.
response := client sendMessage: request.
{ response rateLimits rateLimitRequestsRemaining.
  response rateLimits rateLimitRequestsReset.
  response rateLimits hasRateLimits }
"=> #(98 '2026-04-22T17:30:00Z' true)  (values vary per response)"
\end{lstlisting}

For the rate-limit exception that 429 responses produce see
\S\ref{sec:errors}; for service-tier interactions with rate limits
see \S\ref{sec:service-tiers}.

\medskip
\noindent{\small\textsf{Status: \emph{Available}.\quad
Anthropic docs: \url{https://docs.claude.com/en/api/rate-limits}.}}

\section{Service Tiers}
\label{sec:service-tiers}

Service tiers let the caller trade throughput against cost and
latency for a given request.  The Messages API exposes two values
on the \texttt{service\_tier} request-body field: \texttt{auto} (the
default --- the scheduler picks the best tier available to the
account, including priority when capacity permits) and
\texttt{standard\_only} (force the standard tier even when priority
capacity is available, typically for cost-predictability).
Asynchronous batch processing is not a \texttt{service\_tier} value;
it is reached via a separate endpoint (\S\ref{sec:batches}) whose
requests are billed at the batch discount regardless of what the
per-request \texttt{service\_tier} field would otherwise select.
Anthropic returns the tier actually applied on the response so
callers can observe which tier the scheduler chose.

In Pharo the tier is set via \texttt{serviceTier:} on
\texttt{ClaudeMessageRequest}, which delegates to the underlying
\texttt{ClaudeRequestOptions}.  The setter takes the wire string
directly --- pass either a raw String or one of the
\texttt{ClaudeServiceTier} catalog constants
(\texttt{ClaudeServiceTier auto},
\texttt{ClaudeServiceTier standardOnly}) which return those exact
strings.  The catalog is a discovery surface, not a typed enum:
the setter still accepts arbitrary Strings so unknown values flow
through to Anthropic if it ever extends the field.  Validation is
deliberately left to the server because client-side enforcement
would lock out forward-compat tiers.  The request serializer emits
the field only when a value is set, so requests that leave the tier
unset fall back to Anthropic's default (equivalent to
\texttt{'auto'}).

The example below demonstrates selecting the default tier
explicitly on a message request and reading the configured value
back.  It uses the \texttt{ClaudeMessageRequest} convenience setter,
which writes through to the request options; the body matches
\texttt{ClaudeMessagingRequestOptionsExample}'s canonical build
shape:

\begin{lstlisting}[language={}]
"Build a request that asks the scheduler to pick the best tier."
| request |
request := ClaudeMessageRequest new
    model: ClaudeModel sonnet46;
    maxTokens: 1024;
    serviceTier: 'auto';  "Let the scheduler pick; 'standard_only' forces standard."
    addUserMessage: 'What is the capital of France?';
    yourself.
request serviceTier
"=> 'auto'"
\end{lstlisting}

The tier that Anthropic actually applied to a response is read off
\texttt{ClaudeUsage>>serviceTier}.  The server echoes the chosen
tier in the response body under \texttt{usage.service\_tier} and
\texttt{ClaudeUsage class >> fromJson:} parses that field into the
\texttt{serviceTier} slot; the value is nil when Anthropic does not
echo the field (older endpoints, or count-tokens replies).  Note
that the response-side vocabulary is distinct from the request-side
\texttt{ClaudeServiceTier} catalog: the request says \emph{which
scheduling policy to use} (\texttt{'auto'} or \texttt{'standard\_only'}
today), while \texttt{usage.service\_tier} reports \emph{which
backend tier was actually applied} (typically \texttt{'standard'},
\texttt{'priority'}, or \texttt{'batch'}).  Treat the
\texttt{usage.service\_tier} value as a raw String for now;
\texttt{ClaudeServiceTier} does not enumerate response-side
identifiers, only request-side policies.  The
broader rate-limit telemetry that the live response carries lives
on \texttt{ClaudeMessage>>responseMetadata} (see
\S\ref{sec:rate-limits}), so callers reading the applied tier and
the rate-limit counters drill into two adjacent objects on the same
response rather than re-parsing headers themselves.

\begin{lstlisting}[language={}]
"Read the tier the scheduler applied to a response."
| client request response |
client := ClaudeSDKExampleSupport resolveClient.
request := ClaudeMessageRequest new
    model: ClaudeModel sonnet46;
    maxTokens: 64;
    serviceTier: ClaudeServiceTier auto;
    addUserMessage: 'ping';
    yourself.
response := client sendMessage: request.
{ response usage serviceTier. ClaudeServiceTier allKnown asArray }
"=> #('standard' #('auto' 'standard_only'))  (applied value varies)"
\end{lstlisting}

For the Batches API that owns asynchronous batch-tier processing
see \S\ref{sec:batches}; for the rate-limit interactions the
\texttt{auto} and \texttt{standard\_only} tiers share see
\S\ref{sec:rate-limits}.

\medskip
\noindent{\small\textsf{Status: \emph{Available}.\quad
Anthropic docs: \url{https://docs.claude.com/en/api/service-tiers}.}}

\section{Client SDKs}
\label{sec:client-sdks}

Anthropic publishes first-party client SDKs in Python, TypeScript,
Java, Ruby, C\#, PHP, and a CLI, and maintains a catalog of
additional community and partner SDKs.  Each SDK wraps the same
Messages API over HTTP and exposes the same response shapes; they
differ in idiom, packaging, and what their host ecosystem lets them
express.  The Pharo SDK this specification describes is a
Punt Labs port into the Pharo Smalltalk ecosystem, not
Anthropic-maintained; it ships in two layers, a Messaging layer
(this document) and an Agent layer (separate specification).

This section is reference material; no runnable example applies.

\medskip
\noindent{\small\textsf{Status: \emph{Available}.\quad
Anthropic docs: \url{https://docs.claude.com/en/api/client-sdks}.}}

\section{Batches API}
\label{sec:batches}

The Batches API accepts up to 100\,000 Messages-API requests in a
single job (256\,MB per batch) and runs them asynchronously against
the \texttt{batch} service tier at a discounted per-token rate.
Callers submit a batch, poll for completion, and then stream the
per-request results as a single download once the job finishes.  The
endpoint covers six operations: create, retrieve, retrieve results,
list, cancel, and delete.  Applications that can tolerate minutes of
latency in exchange for the batch-tier discount use this surface to
run overnight evaluations, bulk classification jobs, or warehouse
backfills against the same Messages API they use interactively.

The Pharo SDK does not implement the Batches API today.  No
\texttt{ClaudeBatch*} classes exist in any \texttt{Claude-Messaging-*}
package, and \texttt{ClaudeClient} exposes no batch selectors.  The
implementation plan when it lands mirrors the Files API shape: a
typed request/response layer in \texttt{Claude-Messaging-Types},
\texttt{ClaudeClient} selectors for the six operations, and a
streaming reader for the results download so large batches do not
materialize in full in VM memory.

This section is reference material; no runnable example applies.

\medskip
\noindent{\small\textsf{Status: \emph{Not yet implemented}.\quad
Anthropic docs: \url{https://docs.claude.com/en/api/creating-message-batches}.}}

\section{Files API}
\label{sec:files}

The Files API uploads a file once and lets later requests reference
it by ID instead of inlining base64 content on every call --- you
save tokens, you save bandwidth, and large PDFs or images stop
bloating every request payload.  The Pharo surface covers six Files
API operations via seven selectors on \texttt{ClaudeClient}
(\texttt{uploadFile:} and \texttt{listFiles} each have a one-arg and
a parameterized form), each mirroring one of Anthropic's
\texttt{client.beta.files.*} endpoints in the Python SDK.  The
endpoints all sit behind the \texttt{files-api-2025-04-14} beta
header, which the SDK injects automatically; callers never set it.

\begin{center}
\small
\begin{tabular}{@{}l l p{5.4cm}@{}}
\toprule
\textbf{Selector} & \textbf{HTTP} & \textbf{Returns} \\
\midrule
\texttt{uploadFile:} & POST \texttt{/v1/files} & \texttt{ClaudeFileMetadata}.  One-arg form infers filename and MIME type from the \texttt{FileReference}. \\
\texttt{uploadFile:filename:mimeType:} & POST \texttt{/v1/files} & \texttt{ClaudeFileMetadata}.  Three-arg form to override either. \\
\texttt{listFiles}, \texttt{listFiles:} & GET \texttt{/v1/files} & \texttt{ClaudeFilePage}.  Pass a \texttt{ClaudeFileListParams} for cursor pagination (\texttt{limit}, \texttt{afterId}, \texttt{beforeId}). \\
\texttt{getFileMetadata:} & GET \texttt{/v1/files/\{id\}} & \texttt{ClaudeFileMetadata} \\
\texttt{downloadFile:} & GET \texttt{/v1/files/\{id\}/content} & A \texttt{ByteArray}.  Only Claude-generated files are downloadable; user uploads report \texttt{downloadable\,=\,false}. \\
\texttt{deleteFile:} & DELETE \texttt{/v1/files/\{id\}} & \texttt{ClaudeDeletedFile} \\
\bottomrule
\end{tabular}
\end{center}

The four typed request and response classes carry the API's data
shape: \texttt{ClaudeFileMetadata} (id, filename, mimeType,
sizeBytes, createdAt, downloadable), \texttt{ClaudeFilePage} (data
array, hasMore, firstId, lastId), \texttt{ClaudeFileListParams}
(limit, afterId, beforeId), and \texttt{ClaudeDeletedFile} (id,
deleted).

The most common usage is upload-once-reference-many: you upload a
PDF or image, get a file ID back, and then attach it to message
content blocks via \texttt{ClaudeDocumentBlock fromFileId:} or the
equivalent on \texttt{ClaudeImageBlock}.

The example below demonstrates the read path against
\texttt{ClaudeMockServer} so it runs in the eval server without
an API key or network connection.  It registers two canned
responses (the file inventory and a single file's metadata),
points a \texttt{ClaudeClient} at the mock's local port, calls
\texttt{listFiles} and \texttt{getFileMetadata:}, and extracts
four fields from the typed response objects.  The same body
lives in \texttt{ClaudeClient}'s class comment as the canonical
doctest:

\begin{lstlisting}[language={}]
| server client meta page |
"Spin up an in-image mock HTTP server so the example needs neither network nor API key."
server := ClaudeMockServer startOn: 0.
[ "Register the canned response for GET /v1/files (the inventory listing)."
  server
    respondTo: '/v1/files'
    with: '{"data":[{"id":"file-abc","type":"file","filename":"notes.txt","mime_type":"text/plain","size_bytes":6,"created_at":"2026-04-21T00:00:00Z","downloadable":false}],"has_more":false,"first_id":"file-abc","last_id":"file-abc"}'.
  "Register the canned response for GET /v1/files/file-abc (single-file metadata)."
  server
    respondTo: '/v1/files/file-abc'
    with: '{"id":"file-abc","type":"file","filename":"notes.txt","mime_type":"text/plain","size_bytes":6,"created_at":"2026-04-21T00:00:00Z","downloadable":false}'.
  "Build a ClaudeClient that talks to the mock instead of api.anthropic.com."
  client := ClaudeClient
    apiKey: 'test-key'
    baseUrl: 'http://localhost:' , server port printString
    timeout: 5.
  "Exercise the two read selectors and pull four fields out of the typed responses."
  page := client listFiles.
  meta := client getFileMetadata: 'file-abc'.
  ^ { page data size. page data first id. meta filename. meta sizeBytes } ]
    ensure: [ server stop ]  "Always stop the mock, even on error."
"=> #(1 'file-abc' 'notes.txt' 6)"
\end{lstlisting}

Real upload, download, and delete usage is mechanically the same
shape --- swap the mock server for a production client
(\texttt{ClaudeSDKExampleSupport resolveClient} pulls the API key from
the OS keyring or \texttt{ANTHROPIC\_API\_KEY} and points at
\texttt{api.anthropic.com}) and swap the verbs.
\texttt{client uploadFile: aFileReference} returns a metadata
object whose \texttt{id} is the file ID;
\texttt{client downloadFile: aFileId} returns a raw
\texttt{ByteArray}; \texttt{client deleteFile: aFileId} returns a
\texttt{ClaudeDeletedFile} whose \texttt{deleted} field reflects
success.

Two implementation details are worth knowing.  Uploads stream from
disk in 64\,KB chunks via
\texttt{ClaudeStreamingMultiPartEntity}, so VM memory usage is
bounded by the chunk size plus framing rather than by the file
size --- gigabyte uploads are safe.  Downloads follow the server's
\texttt{downloadable} flag: only Claude-generated files (tool
outputs, skill artifacts) come back; your own uploads are listable
and referenceable but not readable.  Unlike the Skills API's hard
25\,MB per-file and 100\,MB total caps, the Files API has no
client-side size limit --- whatever the server accepts goes
through.

For attaching uploaded files to message content see
\S\ref{sec:content-blocks}; for the sibling Skills API that shares
the streaming multipart machinery see \S\ref{sec:skills}; for the
exception hierarchy that every Files call can raise see
\S\ref{sec:errors}.

\medskip
\noindent{\small\textsf{Status: \emph{Available}.\quad
Anthropic docs: \url{https://docs.claude.com/en/api/files-create}.\quad
Python SDK peer:
\url{https://docs.anthropic.com/en/api/files-list}.}}

\section{Skills API}
\label{sec:skills}

The Skills API packages a directory tree as a reusable, versioned
artifact the model can load into a code-execution container on a
Messages request.  A skill is a folder rooted at \texttt{SKILL.md}
plus whatever supporting files the skill's scripts reference;
uploading that folder gives back a skill id and a version id the
model can mount by reference on any future turn.  The endpoint ships
behind the \texttt{skills-2025-10-02} beta header; every skill URL
also appends \texttt{?beta=true} inline, matching Anthropic's
hosted-beta path convention (pagination requests follow the query
string with \texttt{\&page=}).

The Pharo surface splits across two packages by role.
Request-side specs that bind into the Messages container ---
\cls{ClaudeContainerParams} (the container shape) and
\cls{ClaudeSkillSpec} (a single skill reference) --- live in
\pkg{Claude-Messaging-Types} alongside \cls{ClaudeRequestOptions},
because that is where a caller reaches for them when building a
request.  Metadata responses, pagination wrappers, list-query
params, and deletion receipts live in
\pkg{Claude-Messaging-Skills}.

\begin{center}
\begin{tabular}{@{}lp{6.5cm}@{}}
\toprule
\textbf{Class} & \textbf{Role} \\
\midrule
\cls{ClaudeSkill} & Skill metadata response (id, display title,
  source, latest version) \\
\cls{ClaudeSkillVersion} & Versioned skill metadata response (id,
  name, description, directory, version) \\
\cls{ClaudeDeletedSkill} & Skill deletion response \\
\cls{ClaudeDeletedSkillVersion} & Version deletion response \\
\cls{ClaudeSkillPage} & Page-token-paginated list of skills \\
\cls{ClaudeSkillVersionPage} & Page-token-paginated list of versions \\
\cls{ClaudeSkillListParams} & List query (page, source filter, limit) \\
\cls{ClaudeSkillVersionListParams} & Version list query (page, limit) \\
\cls{ClaudeContainerParams} & Container request shape (id + skills
  array, in \pkg{Claude-Messaging-Types}) \\
\cls{ClaudeSkillSpec} & Skill ref for container loading
  (\texttt{skill\_id} + \texttt{type} + \texttt{version}, in
  \pkg{Claude-Messaging-Types}) \\
\bottomrule
\end{tabular}
\end{center}

\begin{sloppypar}
Both \cls{ClaudeContainerParams} and \cls{ClaudeSkillSpec} carry an
\sel{extraFields} slot so a \texttt{fromJson:}/\texttt{asJson}
roundtrip preserves any wire keys Anthropic ships in the future;
this mirrors the forward-compat pattern \cls{ClaudeMCPServerDefinition}
established for the MCP connector, and means a consumer can read
today's shape today without losing tomorrow's fields.  The beta
header itself is hardcoded on the Skills HTTP client ---
\sel{newSkillHttpClient} emits \texttt{anthropic-beta:
skills-2025-10-02} unconditionally, so callers do not opt in via
\texttt{betas:}, the same shape \sel{newFileHttpClient} uses for
the Files API.
\end{sloppypar}

\begin{sloppypar}
Uploads stream from disk in 64\,KB chunks via
\cls{ClaudeStreamingMultiPartEntity}; even a skill at the 100\,MB
total limit never materializes in VM memory.  The multipart body
uses multiple \texttt{files[]} parts (the literal bracketed field
name is what the server expects --- an unbracketed \texttt{files}
field name is rejected), one part per file under the directory.
The shared helper \sel{buildStreamingSkillPartsInto:fromDirectory:}
assembles those parts for both \sel{uploadDirectory:displayTitle:}
(new skill) and \sel{uploadVersionFor:fromDirectory:} (new version
of an existing skill).  By contrast,
\sel{uploadFile:filename:mimeType:} (Files API,
\S\ref{sec:files}) drives \cls{ClaudeStreamingMultiPartEntity}
directly without a shared helper --- the two APIs diverge enough in
shape that a cross-API helper would cost more than it saves.
Filenames in the multipart parts are prefixed with the upload
directory's basename and use POSIX (forward-slash) separators, so a
file at \texttt{skillDir/SKILL.md} serializes as
\texttt{<basename>/SKILL.md} and a nested
\texttt{skillDir/helpers/util.py} as
\texttt{<basename>/helpers/util.py}; the server requires every
uploaded file to live under a single top-level folder named after
the skill.
\end{sloppypar}

Skills list endpoints return page-token pagination (\texttt{?page=}
request, \texttt{next\_page} response field) rather than the Files
API's cursor-id shape.  \cls{ClaudeSkillPage} and
\cls{ClaudeSkillVersionPage} both expose a \sel{nextPage} slot
(String or nil) and a derived \sel{hasMore} predicate, so a pager
loop reads as a \texttt{[ page hasMore ] whileTrue:} driving the
next \sel{page:} on the list-params.

\begin{sloppypar}
On the container side, \cls{ClaudeRequestOptions} accepts either
shape through its polymorphic \sel{container:} setter --- a plain
String container id (back-compat with the original Code Execution
Tool) or a \cls{ClaudeContainerParams} instance holding an id plus
a skills array.  The dispatch happens in \sel{addFieldsTo:}, which
emits the String verbatim when \texttt{isString} holds and calls
\texttt{asJson} on the params otherwise; existing String-input
callers keep working unchanged.  On the response side,
\cls{ClaudeExecutionContainer} carries a \sel{skills} slot whose
entries parse back as \cls{ClaudeSkillSpec} instances --- the
in-container response shape is identical to the request-side spec
(\texttt{skill\_id}, \texttt{type}, \texttt{version}), so one class
serves both.
\end{sloppypar}

\begin{sloppypar}
Before a single byte leaves the client,
\sel{uploadDirectory:displayTitle:} (via
\sel{buildStreamingSkillPartsInto:fromDirectory:}) enforces three
preconditions and raises a plain \cls{Error} with a pinned
message format for each, so a caller sees a clear precondition
failure instead of a 400 round-trip.  The guards fire in a fixed
order: first, \texttt{SKILL.md} must exist at the upload directory
root; second, the directory basename must match the \texttt{name:}
field in \texttt{SKILL.md}'s YAML frontmatter, which the guard
parses line-by-line inside the leading \texttt{---}~\ldots~\texttt{---}
block (no YAML parser dependency); third, every file must be under
\cls{ClaudeClient} \sel{maxSkillFileBytes} (25\,MB = 26,214,400
bytes) and the summed upload under \sel{maxSkillUploadBytes}
(100\,MB = 104,857,600 bytes).  A skill whose directory name
disagrees with the frontmatter AND whose files are oversized
reports the name mismatch first; the precedence matters because
the caller usually wants to know about the cheaper fix before the
more expensive one.
\end{sloppypar}

Deletion is two-phase by server contract: \sel{deleteSkill:}
rejects a skill that still has versions, so the caller must
enumerate via \sel{listVersionsFor:} and call \sel{deleteVersion:for:}
on each before the skill itself can go.  The SDK does not ship a
cascade helper in v1.0; a follow-up bead tracks adding one if the
two-step dance turns out to be a papercut in practice.

\begin{lstlisting}[language={},title={Skills API}]
"Upload a skill directory; returns a ClaudeSkill metadata response"
| client skill request |
client := ClaudeClient apiKey: 'sk-ant-...'.
skill := client
    uploadDirectory: 'skills/python-data-analysis' asFileReference
    displayTitle: 'Python data analysis'.

"Build a Messages request that loads the skill into its container"
request := ClaudeMessageRequest new
    model: ClaudeModel sonnet46;
    maxTokens: 1024;
    "Polymorphic container: either a String id or ClaudeContainerParams"
    container: (ClaudeContainerParams new
        skills: { ClaudeSkillSpec new
            skillId: skill id;
            type: 'custom';
            "Explicitly pin the current published version; omit version: to let Anthropic use the latest published version"
            version: skill latestVersion;
            yourself };
        yourself);
    addUserMessage: 'Use the python-data-analysis skill to plot ...';
    yourself.
\end{lstlisting}

For the sibling Files API that shares \cls{ClaudeStreamingMultiPartEntity}
see \S\ref{sec:files}; for the feature-guide companion that frames
when to reach for a skill versus a tool see \S\ref{sec:agent-skills};
for the Messages request the container binds into see
\S\ref{sec:messages}.

\medskip
\noindent{\small\textsf{Status: \emph{Available}.\quad
Anthropic docs: \url{https://docs.claude.com/en/api/skills-guide}.}}

\section{Prompt Tools API}
\label{sec:prompt-tools}

Prompt Tools is Anthropic's set of helper endpoints for developing
and improving prompts programmatically: \texttt{generate} drafts a
first-pass prompt from a task description, \texttt{templatize}
extracts variables from a concrete prompt, and \texttt{improve}
proposes revisions to an existing prompt.  The intended use case is
prompt-engineering automation --- CI-driven A/B workflows, batch
prompt iteration, agent build pipelines that evolve their own
prompts --- rather than per-turn production traffic.  The endpoints
live under \texttt{/v1/experimental/}, which is Anthropic's signal
that the shape may change before graduating to stable.

The Pharo SDK does not ship Prompt Tools in v1.0.  The experimental
path is the reason: a typed Smalltalk surface built against
\texttt{/v1/experimental/} today would likely require a breaking
revision when Anthropic moves the endpoints out of that path, and
the SDK's v1.0 promise is to freeze the public types we ship.
\cls{ClaudeClient} exposes no \texttt{generatePrompt:},
\texttt{templatizePrompt:}, or \texttt{improvePrompt:} selectors,
and no request/response classes exist for the three endpoints.
The Phase 2 plan is to add the typed surface when the endpoints
leave \texttt{/v1/experimental/}; until then, teams that need these
capabilities call the endpoints directly via \texttt{ZnClient}
against the raw URL.  This section is reference material; no
runnable example applies.

For the Messages API that the generated and improved prompts
ultimately feed into see \S\ref{sec:messages}.

\medskip
\noindent{\small\textsf{Status: \emph{Not yet implemented}.\quad
Anthropic docs: \url{https://docs.claude.com/en/api/prompt-tools-generate}.}}

\section{Admin API}
\label{sec:admin}

The Admin API is Anthropic's operator-facing control plane: it
manages organizations, workspaces, members, and API keys; it
reports Usage and Cost; it exposes Claude Code Analytics.  The
audience is a platform operator auditing spend, provisioning keys
for new teams, attributing cost across workspaces, or pulling
analytics into a monitoring stack.  Requests are gated by a
distinct \emph{Admin API key}, separate from the standard API key
the Messages surface uses, and minted from the admin console
rather than the developer console.

The Pharo SDK does not ship the Admin API in v1.0.  The scope
declaration in \S\ref{sec:overview} targets the standard messaging
surface --- send, stream, tool use, files, skills --- because that
is what an agent developer reaches for; the Admin surface is a
different product with a different key and a different user.
\cls{ClaudeClient} exposes no admin selectors, and no
organization/workspace/member/key-management classes exist in the
Messaging-layer packages.  The Phase 2 plan, when user demand
materializes, is likely a separate \texttt{ClaudeAdminClient} (or
an opt-in admin surface activated with the Admin API key) rather
than overloading \cls{ClaudeClient}; the two key types and the two
user journeys are distinct enough that blending them into one
client would be more confusing than helpful.  This is operator
surface, not developer surface.  This section is reference
material; no runnable example applies.

For the scope declaration that tracks SDK-level omissions see
\S\ref{sec:overview}.

\medskip
\noindent{\small\textsf{Status: \emph{Not yet implemented}.\quad
Anthropic docs: \url{https://docs.claude.com/en/api/administration-api}.}}


\part{Build with Claude}

\section{Streaming Messages}
\label{sec:streaming}

Streaming delivers a response token-by-token as the model
generates it, rather than forcing the caller to wait for the
full completion.  The transport is Server-Sent Events: the same
\texttt{POST /v1/messages} endpoint as
\S\ref{sec:messages}, just with \texttt{stream: true} on the
request.  Applications reach for streaming whenever a user is
watching the response (chat UIs, IDE assistants, dashboards)
so partial output shows up immediately instead of after a
several-second round-trip.

The Pharo SDK exposes two modes.  The block-callback form
\sel{streamMessage:do:} is the idiomatic path for most
applications: it opens the stream, reads events into typed
dictionaries, hands each event to the caller's block, and
returns the fully accumulated \cls{ClaudeMessage} when the
stream closes.  The raw-socket form \sel{openStreamingSocket:}
is the lower-level path for callers that need to interleave
reads with other work (a \cls{SharedQueue} producer/consumer
pattern, a forked UI updater, etc.); it returns the socket for
the caller to drive directly through \cls{ClaudeSSEDecoder}.
The runnable example below uses \sel{streamMessage:do:} and is
the primary shape a first-time reader should try.  It is the
canonical doctest from \texttt{ClaudeClient}'s class comment:

\begin{lstlisting}[language={}]
| client request final |
client := ClaudeSDKExampleSupport resolveClient.
request := ClaudeMessageRequest new
	model: ClaudeModel opus46;
	maxTokens: 1024;
	addUserMessage: 'Count to five.';
	yourself.
final := client streamMessage: request do: [ :event |
	(event at: 'event') = 'content_block_delta' ifTrue: [
		| delta |
		delta := (event at: 'data') at: 'delta'.
		Transcript show: (delta at: 'text' ifAbsent: [ '' ]) ] ].
Transcript cr; show: 'stop reason: ', final stopReason.
\end{lstlisting}

Each \texttt{event} the block receives is a Dictionary with two
String keys: \texttt{'event'} is the SSE event-type String
(\texttt{'message\_start'}, \texttt{'content\_block\_delta'},
\ldots), and \texttt{'data'} is the parsed-JSON Dictionary
payload.  The block runs synchronously inside
\sel{streamMessage:do:}; when the stream closes,
\sel{streamMessage:do:} returns the accumulated
\cls{ClaudeMessage} with every content block finalized, the
stop reason set, and the token-usage counts populated --- the
same shape a non-streaming \sel{sendMessage:} would have
returned.

\subsection{SSE Protocol}

The streaming API returns Server-Sent Events (SSE) --- a text
protocol where each event is a pair of lines (event type, JSON
data) separated from the next event by a blank line.  Event
types identify what happened (message started, block appended,
delta arrived); the data payload carries the structured
details.  A minimal SSE framing looks like:

\begin{lstlisting}[language={}]
event: message_start
data: {"type":"message_start","message":{...}}

event: content_block_delta
data: {"type":"content_block_delta","delta":{"type":"text_delta","text":"Hello"}}

event: message_stop
data: {"type":"message_stop"}
\end{lstlisting}

\subsection{Event Types}

The eight event types the SDK recognizes map one-to-one to the
steps \cls{ClaudeMessage} walks while assembling a response
from deltas.  \texttt{message\_start} delivers the empty
envelope (id, role, model, initial usage counts);
\texttt{content\_block\_start} announces a new content block
(text, tool use, thinking, \ldots) with its index;
\texttt{content\_block\_delta} carries the incremental payload
keyed on the delta subtype; \texttt{content\_block\_stop}
finalizes the block so streaming accumulators can parse any
partial-JSON buffers; \texttt{message\_delta} updates stop
reason, stop sequence, and the final usage counters;
\texttt{message\_stop} closes the stream; \texttt{ping} is a
heartbeat the SDK ignores silently; \texttt{error} carries an
API error that the client translates into a typed
\cls{ClaudeApiError} subclass:

\begin{center}
\small
\begin{tabular}{@{}lp{7.5cm}@{}}
\toprule
\textbf{Event} & \textbf{Payload} \\
\midrule
\texttt{message\_start}        & Envelope for an empty \cls{ClaudeMessage} (id, role, model, initial usage) \\
\texttt{content\_block\_start} & Content block type and index; creates the right \cls{ClaudeContentBlock} subclass \\
\texttt{content\_block\_delta} & Incremental text, input JSON, thinking, signature, or citation \\
\texttt{content\_block\_stop}  & Block index (finalization signal) \\
\texttt{message\_delta}        & Stop reason, stop sequence, final output-token count \\
\texttt{message\_stop}         & End of stream \\
\texttt{ping}                  & Heartbeat (ignored) \\
\texttt{error}                 & Error during streaming; raised via \texttt{ClaudeApiError fromSseErrorData:} \\
\bottomrule
\end{tabular}
\end{center}

The \texttt{content\_block\_delta} payload is polymorphic,
discriminated by its own \texttt{delta.type} field.  Five delta
subtypes cover every incremental update the server emits; each
maps to a specific field on the active content block:

\begin{center}
\small
\begin{tabular}{@{}llp{6cm}@{}}
\toprule
\textbf{Delta Type} & \textbf{Field} & \textbf{Description} \\
\midrule
\texttt{text\_delta}        & \texttt{text}          & Incremental text content \\
\texttt{input\_json\_delta} & \texttt{partial\_json} & Incremental tool-input JSON (parsed at \texttt{content\_block\_stop}) \\
\texttt{thinking\_delta}    & \texttt{thinking}      & Incremental extended-thinking text \\
\texttt{signature\_delta}   & \texttt{signature}     & Thinking block signature \\
\texttt{citations\_delta}   & \texttt{citation}      & Citation reference appended to the active text block \\
\bottomrule
\end{tabular}
\end{center}

\subsection{SSE Decoder}

\cls{ClaudeSSEDecoder} parses raw SSE text into typed event
dictionaries and supports two calling shapes matching the two
modes above.  The buffered shape --- \sel{parse:do:} --- takes a
complete SSE response body as a String and calls the caller's
two-argument block once per event with
\texttt{(eventType, dataDictionary)}.  Because the full body is
received before the first callback fires, buffered mode is
simpler but loses the latency win.  Callers that already have
the raw body (for example, after \sel{streamRawBody:}) use it
when they just want to iterate the event stream post hoc:

\begin{lstlisting}[language={}]
| body |
"Fetch the full SSE body as a String, then iterate the events."
body := client streamRawBody: request.
ClaudeSSEDecoder new
    parse: body
    do: [ :eventType :data |
        "eventType is a String; data is a Dictionary (or nil on unparseable)."
        self handleEvent: eventType data: data ].
\end{lstlisting}

The incremental shape --- \sel{onStream:} plus \sel{hasNext}
and \sel{nextEvent} --- reads lines off an underlying stream
until a blank line completes an event, then returns the
assembled Dictionary.  This is the shape a \cls{SharedQueue}
producer uses so a consumer process can handle events as they
arrive while the producer reads the socket.  The standard
producer/consumer pattern pairs \sel{openStreamingSocket:}
with \sel{ClaudeSSEDecoder onStream:} and a \cls{SharedQueue};
the producer forks at \texttt{userBackgroundPriority}, the
consumer defers UI updates to the Morphic thread via
\texttt{WorldState defer:}:

\begin{lstlisting}[language={}]
| queue process message |
queue := SharedQueue new.
message := ClaudeMessage new.

"Producer: read the TLS socket in a background process."
process := [ | socket decoder |
  socket := client openStreamingSocket: request.
  decoder := ClaudeSSEDecoder onStream: socket.
  [ decoder hasNext ] whileTrue: [
    queue nextPut: decoder nextEvent ].
  queue nextPut: #end ]
    forkAt: Processor userBackgroundPriority
    named: 'Claude SSE reader'.

"Consumer: process events as they land on the queue."
[ | event |
  event := queue next.
  event == #end ] whileFalse: [
    message accumulate: event.
    "Push partial content to the UI on the Morphic thread."
    WorldState defer: [
      textMorph contents: message textContent ] ].
\end{lstlisting}

\sel{openStreamingSocket:} returns a \cls{ClaudeStreamingSocket}
when the client's \texttt{baseUrl} is an HTTPS endpoint: it
opens a TCP socket, wraps it in \cls{ZdcSecureSocketStream}
with Server Name Indication, sends the HTTP/1.1 POST request
(stream mode true, every required header set), reads the
response status line and headers, and verifies a 2xx status ---
non-2xx raises a typed \cls{ClaudeApiError}.  When the
\texttt{baseUrl} is an HTTP URL (mock servers in tests), the
same selector returns a \cls{ReadStream} over the buffered
body so tests don't need TLS infrastructure.  Both return
types are compatible with the decoder's
\sel{atEnd} / \sel{nextLine} protocol.  Per-read timeout on
the TLS socket comes from \texttt{streamingTimeout} on
\cls{ClaudeClient}, default 300s.  The decoder itself exposes a
small protocol for callers that drive it directly:

\begin{center}
\small
\begin{tabular}{@{}llp{6cm}@{}}
\toprule
\textbf{Method} & \textbf{Returns} & \textbf{Behavior} \\
\midrule
\sel{onStream:}  & \cls{ClaudeSSEDecoder} &
  Class-side factory.  Wraps any object answering
  \sel{atEnd} / \sel{nextLine}.  Initializes the line buffer
  and event accumulator. \\
\sel{hasNext}    & Boolean &
  Reads lines from the stream until a complete event is
  assembled (blank line) or the stream ends.  Returns
  \texttt{true} when an event is ready, \texttt{false} at
  end-of-stream.  Blocks on the socket read --- must run in a
  background Process when the stream is a live socket. \\
\sel{nextEvent}  & Dictionary &
  Returns the last assembled event as
  \texttt{Dictionary('event' -> type. 'data' -> dict)}.
  Call only after \sel{hasNext} returned \texttt{true}. \\
\sel{close}      & --- &
  Closes the underlying stream. \\
\bottomrule
\end{tabular}
\end{center}

Line parsing follows the SSE specification: a line beginning
with \texttt{event:} sets the event type for the pending event,
a line beginning with \texttt{data:} appends to the data
buffer, a blank line dispatches, and a line beginning with
\texttt{:} is a comment and is ignored.
\sel{streamMessage:do:} wraps the decoder loop in a
\texttt{ConnectionTimedOut} handler that re-raises as
\cls{ClaudeTimeoutError} (HTTP 408, retryable at the caller
level) and a \texttt{NetworkError} handler that re-raises as
\cls{ClaudeNetworkError} (synthetic 599, non-retryable);
callers that drive the decoder directly through
\sel{openStreamingSocket:} handle the same two exception
classes themselves.  To cancel an in-flight stream, terminate
the producer process --- the socket is closed on process
termination:

\begin{lstlisting}[language={}]
"Kill the background reader; the socket is closed during termination."
process terminate.
\end{lstlisting}

\subsection{Message Accumulator}

\cls{ClaudeMessage>>accumulate:} builds a complete message
from a sequence of streaming events and is what
\sel{streamMessage:do:} calls under the hood to assemble the
return value.  The method reads the \texttt{'event'} key off
each event dictionary and dispatches to a private handler per
event type: \texttt{message\_start} initializes id, type, turn,
role, model, and usage; \texttt{content\_block\_start}
instantiates the matching \cls{ClaudeContentBlock} subclass
via the type registry and appends it to the active turn's
content; \texttt{content\_block\_delta} sends \sel{applyDelta:}
to the last block, which dispatches polymorphically on the
delta subtype (\texttt{text\_delta} concatenates into
\texttt{text}, \texttt{input\_json\_delta} buffers into
\texttt{partial\_json}, \texttt{thinking\_delta} and
\texttt{signature\_delta} concatenate into
\texttt{thinking}/\texttt{signature}, \texttt{citations\_delta}
appends one \cls{ClaudeCitation} to the active text block's
citations); \texttt{content\_block\_stop} calls
\sel{finalizeAfterStreaming} on the last block so
tool-use/MCP-tool-use blocks can parse their accumulated
partial-JSON buffer into a proper Dictionary;
\texttt{message\_delta} merges the final stop reason, stop
sequence, and updated output-token count; \texttt{error}
raises the matching \cls{ClaudeApiError} subclass via
\texttt{ClaudeApiError class>>fromSseErrorData:}.

The net effect is that after \sel{streamMessage:do:} returns,
the accumulated \cls{ClaudeMessage} is indistinguishable from
the \sel{sendMessage:} response you would have gotten had you
called the non-streaming variant --- same content blocks, same
stop reason, same usage counts.  The Pharo-specific threading
model for streaming (forking the producer process, deferring
to the Morphic thread) lives in the companion
\emph{Pharo-Specific Implementation Notes} under the Pharo
Threading Model section.

For the non-streaming \sel{sendMessage:} counterpart and the
full Messages API surface see \S\ref{sec:messages}; for how
mid-stream SSE errors are routed into the typed exception
hierarchy see \S\ref{sec:errors}; for extended thinking ---
which delivers its own \texttt{thinking\_delta} and
\texttt{signature\_delta} events alongside the normal text
deltas --- see \S\ref{sec:thinking}.

\medskip
\noindent{\small\textsf{Status: \emph{Available}.\quad
Anthropic docs:
\url{https://docs.claude.com/en/docs/build-with-claude/streaming}.\quad
Python SDK peer:
\url{https://docs.anthropic.com/en/api/messages-streaming}.}}

\section{Prompt Caching}
\label{sec:prompt-caching}

Mark a text block, tool definition, or tool-result block as
cacheable, and Anthropic caches the serialized prefix across
requests sharing the same API key --- subsequent requests pay a
fraction of the input-token cost on whatever's cached.  The Pharo
primitive is \texttt{ClaudeCacheControl}, a small value object
that serializes to \texttt{\{"type":"ephemeral"\}} with an
optional \texttt{"ttl"} of \texttt{'5m'} or \texttt{'1h'}.  The
on-wire shape is identical to the Python SDK's
\texttt{cache\_control=\{"type": "ephemeral"\}} block parameter;
the Pharo class just wraps it in a typed object you can pass
around and inspect.

The class-side factories build a marker; \texttt{asJson} converts
it to the wire-format Dictionary.  The default is the
five-minute ephemeral cache; \texttt{ephemeralWithTTL:} sets a
custom TTL (Anthropic currently accepts only \texttt{'5m'} and
\texttt{'1h'}).  Both forms appear below as their own runnable
expressions:

\begin{lstlisting}[language={}]
"Build a default-TTL marker and serialize it for the wire."
ClaudeCacheControl ephemeral asJson
"=> an OrderedDictionary('type'->'ephemeral')"

"Build a marker with an explicit five-minute TTL."
(ClaudeCacheControl ephemeralWithTTL: '5m') asJson
"=> an OrderedDictionary('type'->'ephemeral' 'ttl'->'5m')"
\end{lstlisting}

To use prompt caching in an application, mark a content block as
cacheable and send the request through \texttt{ClaudeClient} as
you would any other.  The most common pattern caches a long
system prompt or RAG prefix.  \texttt{ClaudeMessageRequest}'s
\texttt{systemBlocks:} setter takes a collection of
\texttt{ClaudeTextBlock} instances, one or more of which can
carry a cache marker:

\begin{lstlisting}[language={}]
| client request response |
client := ClaudeClient apiKey: 'sk-ant-...'.

request := ClaudeMessageRequest new
    model: ClaudeModel sonnet46;
    maxTokens: 1024;
    systemBlocks: {
        (ClaudeTextBlock text: 'You are a helpful assistant. <long reusable instructions...>')
            cacheControl: ClaudeCacheControl ephemeral };
    addUserMessage: 'What is the capital of France?';
    yourself.

response := client sendMessage: request.
response usage cacheCreationInputTokens.  "tokens written to cache (first-call miss)"
response usage cacheReadInputTokens.      "tokens served from cache (subsequent hits)"
\end{lstlisting}

\begin{sloppypar}
The first call pays full input cost on the marked prefix and
writes it to the cache (\texttt{cacheCreationInputTokens});
subsequent calls within the TTL pay a fraction
(\texttt{cacheReadInputTokens} reflects what the cache served).
Both counters live on \texttt{ClaudeUsage}, which is on every
Messages response --- no extra API call is needed to observe
cache effectiveness.
\end{sloppypar}

Caching tool definitions and tool-result prefixes follows the
same shape, but the wiring details differ.  The SDK exposes six
\texttt{cacheControl:} setters across six classes
(\texttt{ClaudeMessageRequest}'s is a convenience that delegates to
\texttt{ClaudeRequestOptions}), split into two groups depending on
\emph{when} the marker gets serialized.  \texttt{ClaudeTextBlock}
and \texttt{ClaudeServerTool} take the \texttt{ClaudeCacheControl}
instance directly --- they call \texttt{asJson} on it themselves at
request-serialization time.  \texttt{ClaudeTool},
\texttt{ClaudeToolResultBlock}, \texttt{ClaudeRequestOptions}, and
\texttt{ClaudeMessageRequest} (which delegates to
\texttt{ClaudeRequestOptions}) take the already-serialized
Dictionary --- you call \texttt{asJson} yourself before passing it
in.  The split is deliberate: the
second group's \texttt{cacheControl:} setters predate the typed
wrapper, and tightening them would break callers that already
pass raw dictionaries.  The Pharo type system doesn't catch the
mismatch, so the rule of thumb is: \texttt{ClaudeTextBlock} and
\texttt{ClaudeServerTool} take an instance; everything else takes
\texttt{... asJson}.

\begin{lstlisting}[language={}]
"Cache a tool definition (raw-Dictionary group: pass ... asJson)"
| tool |
tool := ClaudeTool
    name: 'analyze'
    description: 'Analyze code'
    inputSchema: ClaudeToolInputSchema new.
tool cacheControl: ClaudeCacheControl ephemeral asJson.

"Cache a tool-result prefix that the model already saw on the previous turn"
(ClaudeToolResultBlock toolUseId: 'tu_123' content: 'result')
    cacheControl: (ClaudeCacheControl ephemeralWithTTL: '5m') asJson.
\end{lstlisting}

Three constraints are worth planning around.  Anthropic enforces
a maximum of four \texttt{cache\_control} markers per request
server-side, so put them on the largest static prefixes first ---
the system prompt, the tool definitions, the RAG context.  TTL
values currently support only \texttt{'5m'} (five minutes) and
\texttt{'1h'} (one hour); the SDK keeps them as strings so a
future addition like \texttt{'24h'} would need no code change.
And the only \texttt{type} Anthropic exposes is
\texttt{ephemeral}; a future \texttt{persistent} type would land
as a new class-side factory
(\texttt{ClaudeCacheControl persistent}) rather than a new
instance variable.

For the content-block classes that carry cache markers see
\S\ref{sec:content-blocks}; for \texttt{ClaudeTool} and
\texttt{ClaudeServerTool} see \S\ref{sec:tool-use}; for where
\texttt{ClaudeUsage}'s cache fields surface on the response
envelope see \S\ref{sec:messages}.

\medskip
\noindent{\small\textsf{Status: \emph{Available}.\quad
Anthropic docs: \url{https://docs.claude.com/en/docs/build-with-claude/prompt-caching}.\quad
Python SDK peer:
\url{https://docs.anthropic.com/en/docs/build-with-claude/prompt-caching}.}}

\section{Extended Thinking}
\label{sec:thinking}

Extended thinking asks Claude to reason before answering and returns
that reasoning as a separate content-block type the application can
display, hide, or audit.  An agent that writes code, solves a
multi-step math problem, or synthesizes a long document benefits from
visible chain-of-thought: the thinking blocks carry the intermediate
steps, and the final \texttt{ClaudeTextBlock} carries the answer.
Applications that surface reasoning to end users (tutoring,
code-review assistants, evaluation harnesses) read the thinking
blocks alongside the answer; applications that only need the answer
discard them.  The feature is opt-in per request and controlled by a
token budget --- small budgets produce no thinking, large budgets
produce detailed reasoning, and the model chooses how much of the
budget to spend.

\begin{sloppypar}
The Pharo surface has three pieces.  \texttt{ClaudeThinkingParams}
configures the request: \texttt{enabled:} takes an Integer budget in
tokens (required), \texttt{adaptive} lets Anthropic size the budget,
and \texttt{disabled} turns thinking off.  \texttt{ClaudeMessageRequest}
accepts the thinking configuration through its \texttt{thinking:}
setter, which takes a Dictionary --- pass the params object's
\texttt{asDictionary} form.  On the response,
\texttt{ClaudeThinkingBlock} subclasses \texttt{ClaudeContentBlock}
with an \texttt{isThinkingBlock} predicate and \texttt{thinking}
accessor; redacted reasoning (which Anthropic occasionally replaces
for safety) arrives as the sibling
\texttt{ClaudeRedactedThinkingBlock} type carrying opaque
\texttt{data} the application must round-trip verbatim without
inspection.
\end{sloppypar}

The runnable below enables thinking with a 1024-token budget, sends
the request, and filters the response content for the thinking
block.  Anthropic requires \texttt{max\_tokens > budget\_tokens}, so
\texttt{maxTokens} is set to 2048 to leave room for both the
reasoning budget and the final answer:

\begin{lstlisting}[language={}]
| client request response thinkingBlock |
client := ClaudeClient apiKey: 'sk-ant-...'.

"Enable thinking with a 1024-token reasoning budget. The factory
 returns a ClaudeThinkingParams instance; asDictionary converts it
 to the wire-format Dictionary that thinking: expects."
request := ClaudeMessageRequest new
    model: ClaudeModel sonnet46;
    maxTokens: 2048;
    thinking: (ClaudeThinkingParams enabled: 1024) asDictionary;
    addUserMessage: 'What is the Monty Hall problem? Show your reasoning, then give the final answer.';
    yourself.

response := client sendMessage: request.

"Response content is an OrderedCollection of ClaudeContentBlock.
 detect:ifNone: picks the first thinking block; nil-guard the lookup
 because the model may skip thinking on simple prompts even when
 enabled."
thinkingBlock := response content
    detect: [ :b | b isThinkingBlock ]
    ifNone: [ nil ].
thinkingBlock ifNotNil: [ :b |
    Transcript show: b thinking; cr ].

"The final answer is still in the text block; textContent joins all
 text blocks in response order."
response textContent
\end{lstlisting}

Three details shape how an application uses the feature.  Budget
tuning ships today: \texttt{budgetTokens} is a real instance
variable that \texttt{enabled:} sets directly, and a budget below
roughly 1024 tokens often yields no thinking block at all because
the model decides the problem does not warrant deliberation.  Larger
budgets give the model more room to reason but cost more
input-equivalent tokens and extend response latency; the idiomatic
starting point is 1024.  The three modes serve different workloads:
\texttt{enabled:} pins a fixed budget (predictable latency and
cost), \texttt{adaptive} lets Anthropic size the budget to the
prompt (simpler to configure), and \texttt{disabled} turns thinking
off when a previously-enabled request template is being reused for a
quick query.  And the adaptive mode accepts an optional display hint
via \texttt{adaptiveWithDisplay:} (\texttt{'summarized'} or
\texttt{'omitted'}) controlling what the UI receives.

Two further pieces shape multi-turn agents that consume thinking.
The first is the signature round-trip.  Every
\texttt{ClaudeThinkingBlock} returned by the API carries an opaque
cryptographic \texttt{signature} the server computes from the
thinking content plus server-side secrets.  When an application
replays a prior assistant turn back to the model --- the standard
shape of multi-turn tool-use loops --- the signature must arrive at
the server byte-for-byte unchanged or the request is rejected.  The
SDK preserves the signature automatically: \texttt{asJson} on a
\texttt{ClaudeThinkingBlock} emits the \texttt{signature} field, and
\texttt{ClaudeMessageRequest>>asJson} carries it through the
serialized assistant turn.  Applications writing thinking-aware
agents get the round-trip for free; the discipline is pinned by
\texttt{ClaudeThinkingBlockTest>>testSignatureRoundTrip} (the wire
path) and
\texttt{ClaudeThinkingBlockTest>>testStonSignatureRoundTrip} (the
persisted-session path used by repository-backed conversations).

The second is interleaved thinking.  By default the model emits at
most one thinking block per response, before the final text answer.
With the \texttt{interleaved-thinking-2025-05-14} beta the model may
emit a thinking block in \emph{every} assistant turn --- including
the turns that wrap individual tool calls.  Tool-use-heavy agents
benefit because the model can reason about each tool result in turn,
choose the next call deliberately, and reason between calls rather
than only at the end of the response.  Whether that intermediate
reasoning reaches the end user is governed separately by
\texttt{adaptiveWithDisplay:} earlier in this section
(\texttt{'summarized'} or \texttt{'omitted'}).
Opt in per request:

\begin{lstlisting}[language={}]
request betas: { ClaudeBetaHeader interleavedThinking }.
\end{lstlisting}

The runnable
\texttt{ClaudeInterleavedThinkingExample class>>buildRequest}
constructs a multi-turn conversation where two prior assistant
turns each carry their own thinking block plus a tool-use block,
with tool results between them.  The request is structural: its
thinking blocks carry placeholder signatures chosen for clarity
in the doctest, and those placeholders must be replaced with real
signatures captured from a prior live response before the request
can be sent successfully to the API.  The example still doubles
as a test fixture: the placeholder signatures round-trip through
the wire serializer, demonstrating exactly the discipline the
previous paragraph describes.

For \texttt{ClaudeMessageRequest}'s other request knobs and the
in-flight request object see \S\ref{sec:messages}; for how thinking
interacts with tool calls see \S\ref{sec:tool-use}.

\medskip
\begin{sloppypar}
\noindent{\small\textsf{Status: \emph{Available}.\quad
Anthropic docs:
\url{https://docs.claude.com/en/docs/build-with-claude/extended-thinking}.}}
\end{sloppypar}

\section{Context Editing}
\label{sec:context-editing}

Context editing is a server-side set of strategies for trimming
long conversations to fit within Claude's context window.  Anthropic
exposes it as a \texttt{context\_management} request parameter with
named strategies: \texttt{clear\_tool\_uses\_20250919} drops
tool-call and tool-result pairs older than a threshold,
\texttt{clear\_thinking\_20251015} drops thinking blocks from
earlier turns.  Agents that run long tool-use loops accumulate
thousands of tokens of obsolete tool output each turn; a multi-day
conversation retaining every thinking block rapidly exhausts the
window.  Context editing lets the server drop the old material
without the client rebuilding the conversation history locally ---
cheaper than shuttling every turn through a summarizer and more
principled than unbounded growth.

The feature is opted in via the \texttt{context-management-2025-06-27}
beta header (see \S\ref{sec:beta-headers}) plus the
\texttt{context\_management} request parameter naming one or more
strategies and their thresholds.  Anthropic's docs describe the
current strategy set; the set expands as new editing passes land.

\begin{sloppypar}
The Pharo SDK does not implement context editing today.
\texttt{ClaudeRequestOptions} carries no
\texttt{context\_management} slot, \texttt{ClaudeMessageRequest}
exposes no setter for it, and the request serializer does not emit
the field.  When implemented, Pharo will expose a typed
\texttt{ClaudeContextManagement} class with enumerated strategies
and a \texttt{contextManagement:} setter on
\texttt{ClaudeRequestOptions} matching the shape of
\texttt{ClaudeThinkingParams} / \texttt{thinking:}.  This section
is reference material; no runnable example applies.
\end{sloppypar}

For the beta header machinery that will carry the opt-in see
\S\ref{sec:beta-headers}; for the scope declaration that tracks
SDK-level omissions see \S\ref{sec:scope}.

\medskip
\begin{sloppypar}
\noindent{\small\textsf{Status: \emph{Not yet implemented}.\quad
Anthropic docs:
\url{https://docs.claude.com/en/docs/build-with-claude/context-editing}.}}
\end{sloppypar}

\section{Vision / Multimodal}
\label{sec:vision}

Vision lets Claude see images.  The user turn carries a content
block whose type is \texttt{image} alongside any text blocks;
Claude processes the image in its visual encoder and reasons about
it the same way it reasons about text.  Applications use vision for
chart and screenshot analysis, document OCR, photograph
description, UI testing, and any workflow where the input arrives
as pixels rather than characters.  Multiple images per turn are
supported, and images mix freely with text blocks in the same
request.

Pharo wires vision through \texttt{ClaudeImageBlock}, which carries
a \texttt{source} Dictionary shaped for the wire.  Three class-side
factories cover the three source types.  \texttt{fromUrl:} points
at a publicly-accessible image URL; Anthropic fetches and processes
the image server-side.  \texttt{base64:mediaType:} carries the
encoded bytes inline; nothing leaves the client beyond the request
itself, which matters for images the application generates or
derives from private data.  \texttt{fromFileId:} references an
image uploaded earlier through the Files API
(\S\ref{sec:files}); the server looks up the bytes by ID,
avoiding a re-upload on every request when the same image is used
across a long conversation.  Two convenience factories build on
\texttt{base64:mediaType:}: \texttt{fromFile:} reads a local image
file and infers the media type from the extension, and
\texttt{fromForm:} captures a Pharo \texttt{Form} (e.g., a screenshot
from \texttt{World imageForm}) as a PNG.

The runnable below uses \texttt{fromUrl:} because it reads
cleanest.  The image block and a text-block prompt are composed
into a single user turn via \texttt{addUserContent:}, which accepts
a collection of \texttt{ClaudeContentBlock} instances in arbitrary
order:

\begin{lstlisting}[language={}]
| client image request response |
client := ClaudeClient apiKey: 'sk-ant-...'.

"Build the image block from a public URL. Anthropic will fetch and
 process the image server-side; no local IO is required."
image := ClaudeImageBlock fromUrl: 'https://upload.wikimedia.org/wikipedia/commons/4/47/PNG_transparency_demonstration_1.png'.

request := ClaudeMessageRequest new
    model: ClaudeModel sonnet46;
    maxTokens: 1024;
    "addUserContent: returns the added ClaudeMessage, not the request.
     The cascade's yourself targets the original receiver (the request),
     so the assignment captures the request object correctly."
    addUserContent: { (ClaudeTextBlock text: 'Describe this image in one sentence.'). image };
    yourself.

response := client sendMessage: request.
response textContent
\end{lstlisting}

The three factories choose based on where the image lives.  Public
URL is simplest when the image already sits on a CDN or public
host, costs no client-side bandwidth, and Anthropic's fetcher
handles the download.  Base64 inline is the right default for
images the application generates on the fly or derives from
private sources the caller does not want to publish to a public
URL --- the bytes travel in the request body and never leave the
client's control boundary.  File-ID source is the right default
for large images reused across many requests: upload once via the
Files API, pay the upload cost once, and reference the ID on every
subsequent request.  Both \texttt{fromURL:} and \texttt{fromUrl:}
are accepted across \texttt{ClaudeImageBlock} and
\texttt{ClaudeDocumentBlock} (\S\ref{sec:pdf}); either casing
builds the same wire shape.

The inline-source factories validate the raw byte count against
Anthropic's 5\,MB per-image cap before encoding, so oversized
inputs fail fast with a \texttt{ClaudeInvalidRequestError} rather
than waiting for the server's 400 response.  URL and file-ID
sources skip client-side validation --- the server validates at
fetch or lookup time.

Supported formats are PNG, JPEG, GIF, and WebP; Anthropic's docs
carry the current per-image size and dimension limits and the
per-request image count cap.

For the content-block wire shape see \S\ref{sec:content-blocks};
for the Files API used by \texttt{fromFileId:} see
\S\ref{sec:files}; for the sibling document block that handles
PDFs see \S\ref{sec:pdf}.

\medskip
\begin{sloppypar}
\noindent{\small\textsf{Status: \emph{Available}.\quad
Anthropic docs:
\url{https://docs.claude.com/en/docs/build-with-claude/vision}.}}
\end{sloppypar}

\section{PDF / Document Block}
\label{sec:pdf}

PDF support lets Claude read a PDF inline the same way vision lets
it read images.  Anthropic's server decodes the PDF, extracts
text, figures, and tables, and feeds the material into the model's
context alongside the rest of the user turn.  Applications use it
for contract review, research-paper Q\&A, financial-report
summarization, and any workflow where the source document is a PDF
rather than markdown or plain text.  Short PDFs ride inline in the
request body; large or frequently-reused PDFs upload once to the
Files API and reference by ID on every subsequent request.
Anthropic's docs cite a rough budget of 1500--3000 input tokens per
page.

\begin{sloppypar}
Pharo represents documents through \texttt{ClaudeDocumentBlock},
whose \texttt{source} slot holds a \texttt{ClaudeDocumentSource}
subclass matching the wire-format source type.  Four class-side
factories cover the four source types.  \texttt{fromFileId:}
references a PDF uploaded earlier through the Files API
(\S\ref{sec:files}) --- the production-idiomatic path for large or
reused PDFs because it avoids re-uploading the bytes on every
request.  \texttt{fromBase64:} carries the PDF bytes inline as a
base64 string --- simplest for one-shot short PDFs where the
upload round-trip would add more overhead than it saves.
\texttt{fromText:} wraps already-extracted plain text in a document
block --- useful when upstream code has done its own PDF parsing
and wants the text treated as a document (for citations; see
\S\ref{sec:citations}) rather than as an inline user message.
\texttt{fromURL:} points at a publicly-accessible PDF;
\texttt{fromUrl:} works as an alias for callers that prefer the
lowercase convention (\S\ref{sec:vision} covers the same
bidirectional alias on \texttt{ClaudeImageBlock}).
\end{sloppypar}

\begin{sloppypar}
Applications that need to estimate request cost before sending a
large PDF can call \texttt{tokensPerPageRange} on
\texttt{ClaudeDocumentBlock}, which returns an \texttt{Interval}
spanning Anthropic's rough 1500--3000 per-page budget.  Multiply
by the page count for a minimum and maximum input-token estimate,
or call \texttt{ClaudeClient>>countTokens:} for an exact count on
the assembled request.
\end{sloppypar}

The runnable below uses \texttt{fromBase64:} because it runs
end-to-end without a prior Files-API upload.  The document block
and a text-block prompt compose into a single user turn the same
way vision's image block does.  For production use with a large or
reused PDF, swap \texttt{fromBase64:} for \texttt{fromFileId:} with
the identifier returned by \texttt{ClaudeClient>>uploadFile:}:

\begin{lstlisting}[language={}]
| client doc request response |
client := ClaudeClient apiKey: 'sk-ant-...'.

"Carry the PDF bytes inline as a base64 string. The fixture below is
 a minimal 1-page PDF (615 bytes raw, 820 base64 chars) lifted byte-
 for-byte from ClaudeMessagingPDFExample>>examplePDFBase64. No beta
 header is required for base64-inline PDFs."
doc := ClaudeDocumentBlock fromBase64: ClaudeMessagingPDFExample examplePDFBase64.

request := ClaudeMessageRequest new
    model: ClaudeModel sonnet46;
    maxTokens: 1024;
    "addUserContent: returns the added message; the cascade's
     yourself targets the request so the assignment captures it."
    addUserContent: { (ClaudeTextBlock text: 'Summarize the contents of this PDF in one sentence.'). doc };
    yourself.

response := client sendMessage: request.
response textContent
\end{lstlisting}

Four source choices map to four setup costs.  \texttt{fromFileId:}
amortizes a single upload across many requests and is the right
default when the same PDF will be referenced repeatedly (multi-turn
Q\&A over a fixed document).  \texttt{fromBase64:} skips the Files
API entirely --- shortest code path, right when the PDF is short
and used once.  \texttt{fromText:} is the escape hatch when the
caller already has extracted text and wants document-aware
handling (which currently matters chiefly for citations;
non-extracted PDFs are far richer because Anthropic's server
extracts figures and tables too).  \texttt{fromURL:} is the
right default when the PDF lives on a public host the model
should fetch directly.  URL-source PDFs require the
\texttt{document-url-2025-04-15} beta header on the request; the
base64 and file-ID paths have no beta dependency.

For the Files API uploader see \S\ref{sec:files}; for the sibling
image block see \S\ref{sec:vision}; for combining PDFs with
citations see \S\ref{sec:citations}.

\medskip
\begin{sloppypar}
\noindent{\small\textsf{Status: \emph{Available}.\quad
Anthropic docs:
\url{https://docs.claude.com/en/docs/build-with-claude/pdf-support}.}}
\end{sloppypar}

\section{Citations}
\label{sec:citations}

Citations give the application character-indexed pointers back to
the source spans Claude used to produce its answer.  When a user
turn carries a citation-enabled document, each assistant text block
on the response arrives with a \texttt{citations} array naming the
document and the start and end character offsets of the cited
excerpt, plus the excerpt text itself.  Applications use citations
for RAG grounding (the UI can link each sentence of the answer back
to its source), audit trails (a reviewer can verify that every
claim traces to the input), and compliance workflows where the
answer must cite its evidence.  Unlike most Claude behaviour,
citation content is uniquely server-computed: the cited excerpt is
extracted from the raw document bytes the client sent, and the
character offsets are assigned by the server's citation pipeline.
A client cannot fabricate a citation.

\begin{sloppypar}
The Pharo surface has two halves.  On the request, the caller
toggles citations on a \texttt{ClaudeDocumentBlock} with the typed
trio \texttt{enableCitations}, \texttt{disableCitations}, and the
\texttt{citationsEnabled} predicate.  Under the hood the enable
slot holds the wire-shape Dictionary
\texttt{\{ \#enabled -> true \} asDictionary} (mirroring
\texttt{"citations": \{"enabled": true\}} on the JSON side); the
typed methods exist so callers do not have to spell that shape by
hand, and \texttt{disableCitations} sets the slot to \texttt{nil}
so \texttt{asJson} omits the key entirely --- the API recognizes
key omission as the opt-out, not an explicit
\texttt{\{ enabled: false \}}.  The raw \texttt{citations:} setter
remains available as an escape hatch for any future wire-shape
additions.  Only documents carrying the enable marker are eligible
for citation; asking about a plain text block produces no
citations even when the answer is quoted verbatim.  On the
response, each \texttt{ClaudeTextBlock} may carry a
\texttt{citations} OrderedCollection of \texttt{ClaudeCitation}
instances; the accessors are \texttt{citedText} (the extracted
excerpt), \texttt{documentIndex} (the zero-based position of the
cited document in the user turn's content array),
\texttt{documentTitle} (the document's \texttt{title} field when
set), \texttt{startCharIndex}, and \texttt{endCharIndex}.
\end{sloppypar}

The runnable below attaches a citations-enabled document to the
user turn, sends the request, and iterates the response text blocks
to surface every citation:

\begin{lstlisting}[language={}]
| client doc request response citedBlocks |
client := ClaudeClient apiKey: 'sk-ant-...'.

"Build the document block, set an optional title (which appears in
 citations as documentTitle), and enable citations. The typed method
 enableCitations sets the underlying slot to the wire-shape Dictionary
 { #enabled -> true } asDictionary; callers who need to emit a
 future wire extension can still reach for the raw citations: setter."
doc := ClaudeDocumentBlock fromText: 'Pharo images save with open sockets. Restoring creates stale half-open descriptors. The preferred recovery path is to rebuild from clean Tonel source.'.
doc title: 'Pharo Image Discipline'.
doc enableCitations.

request := ClaudeMessageRequest new
    model: ClaudeModel sonnet46;
    maxTokens: 1024;
    "addUserContent: returns the added message; yourself targets the
     request receiver so the assignment captures it."
    addUserContent: { (ClaudeTextBlock text: 'What is the recovery path for a saved image with stale sockets?'). doc };
    yourself.

response := client sendMessage: request.

"Filter response text blocks to the ones carrying citations and
 surface each citation with its excerpt and character range."
citedBlocks := response content select: [ :b |
    b isTextBlock and: [ b citations notNil and: [ b citations isNotEmpty ] ] ].
citedBlocks do: [ :b |
    b citations do: [ :c |
        Transcript
            show: '[', (c documentTitle ifNil: [ '?' ]), ' ',
                c startCharIndex printString, '-', c endCharIndex printString,
                ']: ', c citedText;
            cr ] ]
\end{lstlisting}

Two design details shape how applications use the feature.  The
enable marker lives on the document block, not the request, because
a user turn can mix citation-enabled documents with non-citable
text content --- per-block granularity lets the caller decide which
sources the model may cite.  And the citation objects on the
response point at character offsets in the original document bytes
the client sent, which is what lets the UI highlight the exact
span: the application holds the authoritative document and the
citation tells it where to look.  Prompts that ask about a specific
quoted claim (``what does the document say about X'') yield richer
citations than open-ended summaries, because the model has a clear
target to cite.

For the document block and its four source factories see
\S\ref{sec:pdf}; for citations on image-derived OCR see
\S\ref{sec:vision}; for uploading the source documents through the
Files API see \S\ref{sec:files}.

\medskip
\begin{sloppypar}
\noindent{\small\textsf{Status: \emph{Available}.\quad
Anthropic docs:
\url{https://docs.claude.com/en/docs/build-with-claude/citations}.}}
\end{sloppypar}

\section{Search Results}
\label{sec:search-results}

Search results are a GA content-block type that Anthropic emits
when a RAG document has been searched: each block wraps an excerpt
the model retrieved and carries its own \texttt{citations} array
tying the excerpt back to character offsets in the source.
Applications use them to expose typed RAG responses where each
retrieved snippet already comes citation-linked --- like §\ref{sec:citations}
Citations, but inlined into a dedicated block type rather than
attached to free-form text.  The right time to reach for search
results is when the application runs a RAG pipeline with automatic
citation wiring; sending raw documents and enabling citations on
them (the §\ref{sec:citations} path) stays the right choice when
the application supplies the documents itself and leaves retrieval
to the model.

\begin{sloppypar}
The Pharo SDK does not implement the search-result content block
today.  No \texttt{ClaudeSearchResultBlock} class exists in any
\texttt{Claude-Messaging-*} package, no \texttt{search\_result}
entry is registered in the \texttt{ClaudeContentBlock} type
registry, and the block never decodes on the response --- an
unrecognised type currently falls through to the raw-JSON
preservation path rather than being materialised as a typed
instance.  The server-tool sibling \texttt{ClaudeWebSearchResultBlock}
is a different class and a different feature: it is the result
block the built-in \texttt{web\_search} tool emits (see
§\ref{sec:server-tools}), whereas the RAG \texttt{search\_result}
block is a first-class content-block type the model may produce
directly.  Phase 2 adds a typed \texttt{ClaudeSearchResultBlock}
class, a \texttt{search\_result} entry in the content-block type
registry, and a worked example demonstrating how the block's
internal \texttt{citations} array surfaces through the response.
\end{sloppypar}

This section is reference material; no runnable example applies.

For the response-side citation wire shape that the feature shares
see §\ref{sec:citations}; for the server-tool web-search sibling
see §\ref{sec:server-tools}.

\medskip
\begin{sloppypar}
\noindent{\small\textsf{Status: \emph{Not yet implemented}.\quad
Anthropic docs: \url{https://docs.claude.com/en/docs/build-with-claude/search-results}.}}
\end{sloppypar}

\section{Structured Outputs}
\label{sec:structured-outputs}

Structured outputs constrain the model's response to validate
against a JSON Schema the caller provides.  The model is guaranteed
--- not merely prompted --- to emit a document whose fields match
the schema's property names, value types, and required keys, which
lets the application parse the response as structured data without
retry loops or validation layers on top.  Applications use the
feature for form filling (extract named fields from a user
utterance into a known record type), information extraction
(normalize one shape into another across heterogeneous inputs),
and API glue (treat the model's output as a typed payload the
next service can consume directly).  Opt in by adding the
\texttt{structured-outputs-2025-11-13} beta to the request and
by supplying a format of shape
\texttt{\{type: "json\_schema", schema: \{...\}\}} --- when both
are present the server enforces the schema during generation.
Anthropic's public documentation names the wire parameter
\texttt{output\_format}, but the Pharo SDK's serializer (and the
request body it produces) nests the format inside an
\texttt{output\_config} object, so the on-the-wire JSON field is
\texttt{output\_config.format} carrying the
\texttt{\{type, schema\}} value.  The two are the same feature
under two names; the Pharo SDK emits the grouped form because the
same \texttt{output\_config} object also carries the reasoning
\texttt{effort} hint described later in this section.

\begin{sloppypar}
Pharo exposes the feature through the \texttt{ClaudeOutputFormat}
hierarchy and two request-side setters on \texttt{ClaudeMessageRequest}:
\texttt{outputFormat:} takes a \texttt{ClaudeOutputFormat} instance
(the public API), and \texttt{effort:} takes the optional
reasoning-effort hint (\texttt{'low'}, \texttt{'medium'},
\texttt{'high'}, \texttt{'xhigh'}, or \texttt{'max'}).  Internally
the request aggregates both into a \texttt{ClaudeOutputConfig}
value object held on the private \texttt{outputConfig} slot;
callers reach the format and effort through the two public setters
rather than assigning the config directly, which keeps the request
API symmetric with other parameter groups (sampling params,
thinking params) that also sit behind flat accessors on the
request.
\end{sloppypar}

\begin{sloppypar}
Two concrete format classes subclass \texttt{ClaudeOutputFormat}.
\texttt{ClaudeTextOutputFormat} is the default plain-text output
and is client-side only: its \texttt{asJson} returns \texttt{nil}
so the serializer omits \texttt{output\_config} from the wire
payload entirely (equivalently, it omits what Anthropic's docs
call the \texttt{output\_format} parameter), and its
\texttt{encoding} (\texttt{\#ascii}, \texttt{\#latin1},
\texttt{\#utf8}) and \texttt{markup} (\texttt{\#plain},
\texttt{\#microdown}, \texttt{\#markdown}, \texttt{\#json},
\texttt{\#html}, \texttt{\#latex}) axes translate into a
system-prompt suffix rather than into a wire field.
\texttt{ClaudeJsonSchemaOutputFormat} is the schema-constrained
variant --- it wraps a \texttt{ClaudeToolInputSchema} (the same
builder used for tool definitions, see §\ref{sec:tool-use}) and
serializes to \texttt{\{"type": "json\_schema", "schema": \{...\}\}}.
Construct it via \texttt{ClaudeJsonSchemaOutputFormat class>>schema:}
directly, or via the convenience factory
\texttt{ClaudeOutputFormat class>>jsonSchema:} which delegates to
the same class.
\end{sloppypar}

The runnable below builds a two-field customer-order schema, hands
it to \texttt{outputFormat:} on the request, and sends the message
through \texttt{ClaudeClient}.  The response text then parses as a
JSON document matching the schema:

\begin{lstlisting}[language={}]
| client schema request response |
client := ClaudeClient apiKey: 'sk-ant-...'.

"Build the JSON schema. ClaudeToolInputSchema is the same builder used
 for tool definitions; addProperty:type:description: registers each
 field and require: marks it as mandatory in the emitted schema."
schema := ClaudeToolInputSchema new
    addProperty: 'name' type: 'string' description: 'Customer name';
    addProperty: 'amount' type: 'number' description: 'Order total';
    require: 'name';
    require: 'amount';
    yourself.

"Construct the request. betas: opts into the structured-outputs beta;
 outputFormat: is the public setter that routes the format through
 the internal ClaudeOutputConfig. The broken outputConfig: setter
 some older snippets reference does not exist --- use outputFormat:."
request := ClaudeMessageRequest new
    model: ClaudeModel opus47;
    maxTokens: 1024;
    betas: #('structured-outputs-2025-11-13');
    outputFormat: (ClaudeJsonSchemaOutputFormat schema: schema);
    addUserMessage: 'Record a $42.50 coffee order for Marie Curie as JSON.';
    yourself.

"Send and parse. textContent is the assistant's JSON payload as a String;
 STONJSON fromString: decodes it into a Dictionary whose 'name' value
 is a String and whose 'amount' value is a Number, matching the schema."
response := client sendMessage: request.
STONJSON fromString: response textContent
\end{lstlisting}

The reasoning-effort hint is orthogonal to the output format.  Set
it via \texttt{request effort: 'high'} on any request that supports
it; the two live in the same \texttt{ClaudeOutputConfig} aggregator
but surface as independent levers on the request.

Three worked schema patterns cover most of the shapes a real
extraction request needs.  The first is the \emph{enum-constrained
string field}, used when a value must come from a fixed vocabulary
--- order status, currency code, severity level.  The builder
method \sel{addProperty:type:description:enum:} takes the same
arguments as the unconstrained variant plus an array of allowed
values, and the model is constrained to emit one of them.  Below,
the \texttt{status} field is restricted to three states; the model
cannot return ``in transit'' or any other off-vocabulary value:

\begin{lstlisting}[language={}]
"Enum-constrained string: status must be one of three fixed values."
| schema |
schema := ClaudeToolInputSchema new
    addProperty: 'orderId' type: 'string' description: 'Order identifier';
    addProperty: 'status'
        type: 'string'
        description: 'Order fulfilment state'
        enum: #('pending' 'shipped' 'delivered');
    require: 'orderId';
    require: 'status';
    yourself.
\end{lstlisting}

The second is the \emph{nested object shape}, used when a record
contains a sub-record --- a customer with an address, an invoice
with a billing party, a shipment with origin and destination.  The
builder method \sel{addProperty:schema:} takes a property name and
another \cls{ClaudeToolInputSchema} and registers the sub-schema as
the property's value.  Each sub-schema is built the same way the
top-level one is, so the pattern composes recursively to whatever
depth the wire schema needs:

\begin{lstlisting}[language={}]
"Nested object: a Customer record with an embedded Address."
| address customer |
address := ClaudeToolInputSchema new
    addProperty: 'street' type: 'string' description: 'Street address';
    addProperty: 'city' type: 'string' description: 'City name';
    addProperty: 'postalCode' type: 'string' description: 'Postal code';
    require: 'street';
    require: 'city';
    yourself.
customer := ClaudeToolInputSchema new
    addProperty: 'name' type: 'string' description: 'Customer name';
    addProperty: 'address' schema: address;
    require: 'name';
    require: 'address';
    yourself.
\end{lstlisting}

The third is the \emph{array-of-record payload}, used when the
response is a homogeneous list of records --- order line items, a
slice of search results, a batch of extracted entities.  Set the
schema's top-level \texttt{type} to \texttt{'array'} via
\sel{type:}, then attach the per-element schema with \sel{items:}.
The resulting wire schema is \texttt{\{type: "array", items:
\{...\}\}} and the model emits a JSON array whose entries each
match the inner schema:

\begin{lstlisting}[language={}]
"Array of records: a list of order line items. Each item must
 carry sku and quantity; price is optional so back-orderable lines
 can be quoted later without failing schema validation."
| lineItem orderLines |
lineItem := ClaudeToolInputSchema new
    addProperty: 'sku' type: 'string' description: 'Stock keeping unit';
    addProperty: 'quantity' type: 'integer' description: 'Units ordered';
    addProperty: 'price' type: 'number' description: 'Unit price (optional for back-orderable lines)';
    require: 'sku';
    require: 'quantity';
    yourself.
orderLines := ClaudeToolInputSchema new
    type: 'array';
    items: lineItem;
    yourself.
\end{lstlisting}

\begin{sloppypar}
The example illustrates the \sel{require:} versus optional
distinction: a property registered with
\sel{addProperty:type:description:} appears in the schema's
\texttt{properties} map, but only properties named via
\sel{require:} are listed in the wire schema's \texttt{required}
array.  The model is free to omit any property not in that array
--- here, \texttt{price} can be absent on a back-orderable line
item without violating the schema.
\end{sloppypar}

\begin{sloppypar}
For arbitrary JSON-Schema keywords the builder methods do not
cover --- \texttt{anyOf} and \texttt{oneOf} union types, format
validators, length constraints, regex patterns ---
\sel{schemaField:put:} stores the keyword on the schema's
\texttt{extraFields} dictionary so it round-trips through
\sel{asJson}.  The setter is selective: the four keys the schema
already models as typed state --- \texttt{type}, \texttt{description},
\texttt{properties}, \texttt{required} --- are routed through their
dedicated setters (\sel{type:}, \sel{description:}, \sel{properties:},
\sel{required:}) so internal state stays consistent; every other
keyword lands in \texttt{extraFields} verbatim.  This is
\cls{ClaudeToolInputSchema}'s escape hatch for the full JSON-Schema
surface: anything the schema does not model as a typed builder method
can be added by name.  For example, \texttt{schema schemaField: 'oneOf'
put: \{ stringSchema asJson. numberSchema asJson \}} attaches a
discriminated union as a top-level keyword, and the wire payload
carries it through unchanged.
\end{sloppypar}

\begin{sloppypar}
For where the request parameter fits on the Messages API payload
see §\ref{sec:messages}; for the shared
\texttt{ClaudeToolInputSchema} used in tool definitions see
§\ref{sec:tool-use}.
\end{sloppypar}

\medskip
\begin{sloppypar}
\noindent{\small\textsf{Status: \emph{Available}.\quad
Anthropic docs: \url{https://docs.claude.com/en/docs/build-with-claude/structured-outputs}.}}
\end{sloppypar}

\section{Token Counting (Feature Guide)}
\label{sec:token-counting}

Token counting before you send a request is a planning surface,
not an API surface.  This section is the feature-guide companion
to §\ref{sec:count-tokens} Count Tokens: §\ref{sec:count-tokens}
documents the \texttt{countTokens:} selector and its response
shape, while this section covers the \emph{why} --- pre-flight cost
quoting, UX guardrails that warn the user before a long turn
overruns the model's input cap, context-window selection when
deciding between a 200K-token model and a 1M-token model, and
batch-job gating where skipping a request that would exceed a
quota is cheaper than letting it fail at send time.  The same
count the billing path will meter is available before the request
leaves the client, so the application can price, gate, and size
the turn in one pre-flight call.

\begin{sloppypar}
To estimate tokens in an application the caller builds the
\texttt{ClaudeMessageRequest} the same way \texttt{sendMessage:}
takes one, hands it to \texttt{countTokens:}, and reads
\texttt{inputTokens} off the returned
\texttt{ClaudeMessageTokensCount}.  For the runnable call and the
precise response shape see §\ref{sec:count-tokens}; the API surface
is reference material and lives there rather than being duplicated
here.  The worksheet below uses the same call shape as the
§\ref{sec:count-tokens} example to drive a four-step cost-quoting
flow.
\end{sloppypar}

The first sizing decision is which context window the request needs.
Three families ship today: 200K is the universal baseline; 500K is
the production-tier window on \cls{ClaudeModel} \sel{opus47} and
\sel{opus46}; 1M is the long-context window on Sonnet 4.6, Opus 4.6,
and Opus 4.7 (which include it at standard pricing without a beta
flag, per Anthropic's pricing page) and on Sonnet 4.5 (which gates
it behind the \texttt{context-1m-2025-08-07} beta).  The table below
maps each window to the model families the SDK ships typed
\cls{ClaudeModel} accessors for, and the beta header (if any) the
request must carry --- callers reach the beta string through the
typed catalog \texttt{ClaudeBetaHeader context1m}, which returns
\texttt{'context-1m-2025-08-07'} (see §\ref{sec:beta-headers}).
Each row lists the model families whose maximum input window reaches
the row's size, and the beta column names the only gate the request
must clear --- the 1M opt-in for Sonnet 4.5.  Larger-cap families
appear in smaller-cap rows too because a 1M-capable model still
serves any request that fits a smaller window; the SDK exposes no
per-request knob to select a window below a model's maximum:

\begin{center}
\small
\begin{tabular}{@{}lp{6.0cm}p{3.6cm}@{}}
\toprule
\textbf{Window} & \textbf{Model families} & \textbf{Beta required} \\
\midrule
200K & Sonnet 4.5, Sonnet 4.6, Opus 4.x, Haiku 3.5, Haiku 4.5 & none \\
500K & Opus 4.6, Opus 4.7 & none \\
1M   & Sonnet 4.6, Opus 4.6, Opus 4.7 & none (standard pricing) \\
1M   & Sonnet 4.5 & \texttt{context-1m-2025-08-07} \\
\bottomrule
\end{tabular}
\end{center}

The second decision is what each input token will cost.  Anthropic
publishes per-model rates per million tokens (MTok) on its pricing
page; the figures below are base input and output rates retrieved
from \url{https://docs.claude.com/en/docs/about-claude/pricing} on
2026-04-22.  Cache-write and cache-hit pricing is a separate axis
covered in §\ref{sec:prompt-caching}.  Callers should re-check the
console before quoting --- Anthropic adjusts rates and the table
below is a snapshot, not a contract:

\begin{center}
\small
\begin{tabular}{@{}lrr@{}}
\toprule
\textbf{Model} & \textbf{Input \$/MTok} & \textbf{Output \$/MTok} \\
\midrule
Haiku 3.5  & \$0.80 & \$4.00 \\
Haiku 4.5  & \$1.00 & \$5.00 \\
Sonnet 4.5 & \$3.00 & \$15.00 \\
Sonnet 4.6 & \$3.00 & \$15.00 \\
Opus 4.6   & \$5.00 & \$25.00 \\
Opus 4.7   & \$5.00 & \$25.00 \\
\bottomrule
\end{tabular}
\end{center}

The pre-flight cost quote runs in four steps.  First, build the
\cls{ClaudeMessageRequest} the way \sel{sendMessage:} takes one ---
same model, same conversation, same system prompt.  Second, hand it
to \sel{countTokens:} on the client; the call returns a
\cls{ClaudeMessageTokensCount} carrying the single
\sel{inputTokens} field, the same value the billing path will meter
when the request is sent.  Third, multiply that count by the model's
per-million input rate from the table above; output cost is bounded
by \texttt{maxTokens}, so a worst-case ceiling multiplies it by the
output rate.  Fourth, surface the figure in the UI before the
caller commits.  The runnable below quotes a Sonnet 4.6 turn with
4096 max output tokens; the call shape mirrors the
§\ref{sec:count-tokens} runnable:

\begin{lstlisting}[language={},title={Pre-flight cost quote}]
| client request count inputRate outputRate inputCost ceiling |
client := ClaudeClient apiKey: 'sk-ant-...'.

"Step 1: build the request the same way sendMessage: takes one."
request := ClaudeMessageRequest new
    model: ClaudeModel sonnet46;
    maxTokens: 4096;
    addUserMessage: 'Summarise the attached transcript in three bullet points.';
    yourself.

"Step 2: count input tokens; same shape as the section 6 runnable."
count := client countTokens: request.
count inputTokens  "=> Integer, e.g. 12384"

"Step 3: multiply by per-model rates from the published pricing table."
inputRate := 3.00.   "Sonnet 4.6 input: $3.00 / MTok (retrieved 2026-04-22)"
outputRate := 15.00. "Sonnet 4.6 output: $15.00 / MTok"
inputCost := count inputTokens * inputRate / 1000000.
ceiling := inputCost + (request maxTokens * outputRate / 1000000).

"Step 4: surface the quote in the UI."
String streamContents: [ :s |
    s nextPutAll: 'Input cost: $';
      nextPutAll: (inputCost printShowingDecimalPlaces: 4);
      nextPutAll: ' (worst case with maxTokens output: $';
      nextPutAll: (ceiling printShowingDecimalPlaces: 4);
      nextPutAll: ')' ]
\end{lstlisting}

Two tactics extend the same flow.  Batch jobs gate at step 2: if
\sel{inputTokens} exceeds the per-job quota the orchestrator skips
the send entirely, paying nothing instead of a 400 round-trip.  UX
guardrails fire between steps 2 and 4: if the count is within 90\%
of the chosen model's context window, the application warns the
user before committing rather than after the model truncates.  Both
tactics use the same \sel{countTokens:} call --- the planning surface
is one endpoint, the policy is application-side.

For the \texttt{countTokens:} API reference and the runnable
example see §\ref{sec:count-tokens}; for the beta-header mechanism
the 1M context-window beta uses see §\ref{sec:api-versions}; for
the prompt-caching pricing axis see §\ref{sec:prompt-caching}.

\medskip
\begin{sloppypar}
\noindent{\small\textsf{Status: \emph{Available}.\quad
Anthropic docs: \url{https://docs.claude.com/en/docs/build-with-claude/token-counting}.}}
\end{sloppypar}

\section{Batch Processing}
\label{sec:batch-processing}

Batch processing is the feature-guide companion to §\ref{sec:batches}
Batches API: §\ref{sec:batches} catalogues the six operations
(create, retrieve, retrieve results, list, cancel, delete) that
would land once the surface ships, while this section covers the
\emph{why} --- overnight evaluation runs against large question
banks, bulk classification jobs over warehouse exports,
model-versus-model A/B comparisons over fixed test sets, and
retrospective backfills where days of historical traffic replay
through a new prompt.  Applications reach for the Batches API
when the workload tolerates up to 24 hours of completion latency
in exchange for the discounted batch service tier, which prices
per-token consumption below the interactive \texttt{auto} and
\texttt{standard\_only} tiers.  The batch tier is itself not a
\texttt{service\_tier} request-field value (see
§\ref{sec:service-tiers}); it is intrinsic to the batch endpoint
and applies to every request inside the job automatically.

The Pharo SDK does not implement the Batches API today.  See
§\ref{sec:batches} for the planned surface and the implementation
shape the Pharo layer would take when the feature lands --- a typed
request/response layer in \texttt{Claude-Messaging-Types},
\texttt{ClaudeClient} selectors for the six operations, and a
streaming reader for the results download so jobs with 100\,000
per-request results do not materialize in full in VM memory.
§\ref{sec:scope} lists the Batches API among the v2 surfaces
deferred beyond the v1 release.

This section is reference material; no runnable example applies.

For the API inventory and the planned implementation shape see
§\ref{sec:batches}; for the distinction between the batch endpoint
and the per-request \texttt{service\_tier} field see
§\ref{sec:service-tiers}.

\medskip
\begin{sloppypar}
\noindent{\small\textsf{Status: \emph{Not yet implemented}.\quad
Anthropic docs: \url{https://docs.claude.com/en/docs/build-with-claude/batch-processing}.}}
\end{sloppypar}


\part{Agents and Tools}

\section{Tool Use}
\label{sec:tool-use}

Tool use turns the Messages API into a two-way conversation between
Claude and application code.  The application declares one or more
tools on the request --- each one a name, a prose description, and a
JSON Schema for its inputs.  When Claude decides it needs the tool,
it returns a \texttt{ClaudeToolUseBlock} in the response content;
the caller runs whatever the tool does, wraps the answer in a
\texttt{ClaudeToolResultBlock}, appends it to the message history,
and sends again.  Claude integrates the result and continues.  The
loop is what powers agentic applications: database lookups, API
calls, file reads, shell commands, anything the model can't do on
its own.

The messaging layer ships the primitives for building one turn of
that loop.  A tool is a \texttt{ClaudeTool} --- three slots (name,
description, input schema), plus an optional \texttt{cacheControl}
marker for caching tool definitions across requests
(\S\ref{sec:prompt-caching}).  The input schema is a
\texttt{ClaudeToolInputSchema}, a builder over the JSON Schema
shape the Messages API expects: type defaults to \texttt{'object'},
a \texttt{properties} dictionary holds the parameters,
\texttt{required} lists the mandatory keys.  A tool is attached to a
request via \texttt{ClaudeMessageRequest>>tools:}, which takes a
collection.  When Claude uses the tool, the matching content block
in the response answers \texttt{isToolUseBlock} with \texttt{true}
and exposes \texttt{name} (the tool it picked) and \texttt{input}
(the arguments it supplies as a Dictionary).  The next turn replies
with a \texttt{ClaudeToolResultBlock} whose \texttt{toolUseId} links
back to the block Claude emitted.

The example below shows one turn end-to-end: declare a
\texttt{get\_weather} tool with a single \texttt{city} string
property, attach it to a request, send, and dispatch on the
response content.  The full round-trip loop --- detect the
\texttt{tool\_use} block, run the tool locally, append a
\texttt{tool\_result} block, send again, repeat until Claude stops
requesting tools --- lives in the Claude Agent SDK via
\texttt{ClaudeExchange}.  This section shows only the messaging
primitives; \texttt{ClaudeExchange} is specified in the
\emph{Claude Agent SDK Specification}.

\begin{lstlisting}[language={}]
| client schema tool request response |
client := ClaudeClient apiKey: 'sk-ant-...'.

"Build a JSON Schema with one string property."
schema := ClaudeToolInputSchema new
    addProperty: 'city' type: 'string' description: 'City name';
    yourself.

"Wrap name, description, and schema in a ClaudeTool."
tool := ClaudeTool
    name: 'get_weather'
    description: 'Get current weather for a city'
    inputSchema: schema.

"Attach the tool to a request via tools: (takes a collection)."
request := ClaudeMessageRequest new
    model: ClaudeModel sonnet46;
    maxTokens: 1024;
    tools: (Array with: tool);
    addUserMessage: 'What is the weather in Paris?';
    yourself.

"Send. Claude may answer with a ClaudeToolUseBlock in the response content."
response := client sendMessage: request.
response content do: [ :block |
    block isToolUseBlock ifTrue: [
        "block name   => 'get_weather' "
        "block input  => a Dictionary('city'->'Paris' ) "
        Transcript show: 'Call ', block name, ' with ', block input printString; cr ] ]
\end{lstlisting}

The snippet above was run against a live Pharo image.  Constructing
\texttt{tool} yields \texttt{'get\_weather'} as its \texttt{name};
attaching it via \texttt{tools:} produces a serialized request
whose \texttt{tools} array has one entry; a synthesized
\texttt{ClaudeToolUseBlock} answers \texttt{isToolUseBlock} with
\texttt{true} and round-trips its \texttt{name} and \texttt{input}
slots, with \texttt{input} printing as a Pharo Dictionary
(\texttt{\#(true 'get\_weather' 'a Dictionary(''city''->''Paris'' )')}).

Phase 2 adds the tool-choice modes (\texttt{any},
\texttt{required}, \texttt{tool}),
\texttt{disableParallelToolUse}, fine-grained tool-input
streaming, and the token-efficient-tools beta identifier.  The
messaging primitives above do not change; the request-shaping
setters land on \texttt{ClaudeMessageRequest} and
\texttt{ClaudeRequestOptions}.

For the Anthropic-hosted tool siblings see \S\ref{sec:server-tools};
for external tools bridged through Anthropic see
\S\ref{sec:mcp-connector}; for the request/response shape the tools
ride on see \S\ref{sec:messages}; for thinking blocks that can
interleave with tool-use turns see \S\ref{sec:thinking}; the full
tool-use loop is specified in the \emph{Claude Agent SDK
Specification} via \texttt{ClaudeExchange}.

\medskip
\begin{sloppypar}
\noindent{\small\textsf{Status: \emph{Available}.\quad
Anthropic docs: \url{https://docs.claude.com/en/docs/agents-and-tools/tool-use/overview}.\quad
Python SDK peer:
\url{https://docs.anthropic.com/en/docs/build-with-claude/tool-use/overview}.}}
\end{sloppypar}

\section{Server Tools}
\label{sec:server-tools}

Server tools flip the execution boundary.  Instead of Claude
asking the application to run a tool and waiting for the answer,
Anthropic hosts the tool and runs it directly during message
generation.  The application declares which hosted tools are
allowed on the request; the API invokes them as needed and returns
the results inline with the rest of the response content.  No
round-trip through application code is needed for web search, web
fetch, sandboxed code execution, sandboxed bash, a server-side
text editor, persistent memory, or BM25/regex tool search.
Server tools are what you reach for when you want Claude to
\emph{do} something the API can perform on your behalf, rather
than delegate the work back to your stack.

\begin{sloppypar}
The Smalltalk SDK models each hosted tool as a concrete subclass
of \texttt{ClaudeServerTool} in \pkg{Claude-Messaging-Tools}.
Each class carries the wire-format name, defaults for any
tool-specific parameters, and an \texttt{asJson} method that
serializes the tool declaration Anthropic expects in the request's
\texttt{tools} array.  The hierarchy is small and flat:
\end{sloppypar}

\begin{center}
\begin{tabular}{@{}llp{5cm}@{}}
\toprule
\textbf{Class} & \textbf{Tool Name} & \textbf{Purpose} \\
\midrule
\texttt{ClaudeServerTool}          & (abstract)                  & Base class \\
\texttt{ClaudeWebSearchTool}       & \texttt{web\_search}        & Web search \\
\texttt{ClaudeWebFetchTool}        & \texttt{web\_fetch}         & Fetch URL content \\
\texttt{ClaudeCodeExecutionTool}   & \texttt{code\_execution}    & Sandboxed code execution \\
\texttt{ClaudeBashTool}            & \texttt{bash}               & Sandboxed bash \\
\texttt{ClaudeTextEditorTool}      & \texttt{text\_editor}       & Server-side text editor \\
\texttt{ClaudeMemoryTool}          & \texttt{memory}             & Persistent memory \\
\texttt{ClaudeToolSearchBm25Tool}  & \texttt{tool\_search\_bm25} & BM25 tool search \\
\texttt{ClaudeToolSearchRegexTool} & \texttt{tool\_search\_regex} & Regex tool search \\
\bottomrule
\end{tabular}
\end{center}

\begin{sloppypar}
\texttt{ClaudeMessageRequest} gains an extension method per tool
in \pkg{Claude-Messaging-Tools} --- \texttt{addWebSearch},
\texttt{addWebFetch}, \texttt{addCodeExecution}, \texttt{addBash},
\texttt{addTextEditor}, \texttt{addMemory},
\texttt{addToolSearchBm25}, \texttt{addToolSearchRegex} --- each
one a one-line wrapper that instantiates the matching subclass
with defaults and appends it to the request's tool list.
Because the extension lives in the tools package, the core
\texttt{Claude-Messaging-Types} package stays free of dependencies
on the server-tool hierarchy; loading
\pkg{Claude-Messaging-Tools} into the image is what lights the
convenience methods up.  Server tools can coexist with
client-defined tools (\S\ref{sec:tool-use}) on the same request:
both land in the same \texttt{tools} array, and Anthropic
dispatches to the right executor by tool name.
\end{sloppypar}

\begin{lstlisting}[language={}]
| client params response |
client := ClaudeClient apiKey: 'sk-ant-...'.

params := ClaudeMessageRequest new
    model: ClaudeModel sonnet46;
    maxTokens: 1024;
    addUserMessage: 'Summarise today''s headline at anthropic.com';
    yourself.

"Four one-liners, each appending a ClaudeServerTool subclass to the tools array."
params addWebSearch.
params addWebFetch.
params addCodeExecution.
params addBash.

"The request now carries four hosted-tool declarations."
(params asJson at: 'tools') size.  "=> 4"

"Send. Anthropic runs the tools server-side and returns their results inline."
response := client sendMessage: params.
\end{lstlisting}

\begin{sloppypar}
Reference docs for each hosted tool live under
\url{https://docs.claude.com/en/docs/agents-and-tools/tool-use/}
(\texttt{web-search-tool}, \texttt{web-fetch-tool},
\texttt{code-execution-tool}, \texttt{bash-tool},
\texttt{text-editor-tool}, plus Memory and Tool Search
sub-pages).  Tool-specific betas --- when Anthropic gates a new
server tool behind an opt-in --- ride on \texttt{betas:} in the
request the same way every other beta does
(\S\ref{sec:beta-headers}).
\end{sloppypar}

Phase 2 evaluates Computer Use
(\texttt{computer\_20250124}) and Anthropic's programmatic
tool-calling shape; neither ships in the current release.

For the client-defined tool sibling see \S\ref{sec:tool-use}; for
external tools bridged through an MCP server see
\S\ref{sec:mcp-connector}; for the opt-in header mechanism see
\S\ref{sec:beta-headers}.

\medskip
\begin{sloppypar}
\noindent{\small\textsf{Status: \emph{Available}.\quad
Anthropic docs: \url{https://docs.claude.com/en/docs/agents-and-tools/tool-use/}.\quad
Python SDK peer:
\url{https://docs.anthropic.com/en/docs/build-with-claude/tool-use/overview}.}}
\end{sloppypar}

\section{MCP Connector}
\label{sec:mcp-connector}

The MCP connector beta lets the SDK point Claude at a Model
Context Protocol server hosted somewhere on the public internet,
and have Anthropic's infrastructure --- not the application's
runtime --- connect out to that server during message generation.
The caller lists one or more server URLs on the request under
\texttt{mcpServers:}, opts into the
\texttt{mcp-client-2025-04-04} beta via \texttt{betas:}, and
sends.  Anthropic plays the role of MCP client; when the model
invokes a tool on one of the listed servers, Anthropic dispatches
the call, collects the result, and folds both the invocation and
the result into the response stream as new content-block types
the SDK dispatches on.  The MCP server itself runs externally, on
the customer's infrastructure or a third-party host.

Three features share the MCP acronym in this codebase, and only
one of them is this connector.  \emph{This} connector is the
SDK-side feature where Claude consumes MCP tools through
Anthropic; exposing the Pharo image \emph{as} an MCP server for
another MCP host to call is a separate, post-v1.0 feature; the
local MCP runtime Claude Code uses to back slash commands like
\texttt{/find} or \texttt{/who} is a different surface entirely
(biff, quarry, lux plugins, not an SDK concern).

\begin{sloppypar}
The Smalltalk surface lives in two packages.  Request-side class
definitions are in \pkg{Claude-Messaging-MCP}; the response-side
content blocks live in \pkg{Claude-Messaging-Types} alongside the
other \texttt{ClaudeContentBlock} subclasses, registered in the
name-keyed \texttt{TypeRegistry} under \texttt{mcp\_tool\_use}
and \texttt{mcp\_tool\_result}:
\end{sloppypar}

\begin{center}
\begin{tabular}{@{}lp{6.5cm}@{}}
\toprule
\textbf{Class} & \textbf{Role} \\
\midrule
\texttt{ClaudeMCPServerDefinition}   & Request-side server entry:
  name, URL, optional auth token, optional tool config \\
\texttt{ClaudeMCPToolConfiguration}  & Per-server tool config
  (enabled flag, allow-listed tool names) \\
\texttt{ClaudeMCPToolUseBlock}       & Response block when Claude
  invokes a tool on the MCP server \\
\texttt{ClaudeMCPToolResultBlock}    & Response block carrying the
  result (String or \texttt{OrderedCollection} of
  \texttt{ClaudeTextBlock}) \\
\bottomrule
\end{tabular}
\end{center}

The \texttt{ClaudeMCPToolResultBlock} \texttt{content} slot is
deliberately polymorphic: a wire \texttt{string} stays a
\texttt{String}; a wire array becomes an
\texttt{OrderedCollection} of \texttt{ClaudeTextBlock}.
\texttt{asJson} dispatches on the runtime shape so a
\texttt{fromJson:/asJson} round-trip preserves wire fidelity in
both directions.  Both \texttt{ClaudeMCPServerDefinition} and
\texttt{ClaudeMCPToolConfiguration} carry an \texttt{extraFields}
slot for forward-compat preservation of unknown wire keys,
mirroring the strategy \texttt{ClaudeContentBlock} uses for
unrecognized block types.

The SDK does not auto-inject the MCP beta header when
\texttt{mcpServers:} is set; the caller opts in explicitly via
\texttt{betas: \#('mcp-client-2025-04-04')}.  Explicit opt-in
matches the shape of the Files API, PDF document blocks, and the
\texttt{output-128k} beta --- the SDK never activates a beta on
the caller's behalf.  Two required-field guards fail fast before
the request leaves the client: \texttt{asJson} on
\texttt{ClaudeMCPServerDefinition} raises an \texttt{Error} (via
\texttt{self error:}) when \texttt{name} or \texttt{url} is
\texttt{nil}, and on \texttt{ClaudeMCPToolResultBlock} when
\texttt{toolUseId} is \texttt{nil}.  Both are required on the
wire; failing in-image is preferable to a \texttt{400} round-trip.

\begin{lstlisting}[language={}]
| client server params response |
client := ClaudeClient apiKey: 'sk-ant-...'.

"Build a server entry with name, URL, and optional bearer token."
server := ClaudeMCPServerDefinition new
    name: 'example';
    url: 'https://mcp.example.com/sse';
    authorizationToken: 'token-123';
    yourself.

"Opt into the MCP beta and attach the server list on the request."
params := ClaudeMessageRequest new
    model: ClaudeModel sonnet46;
    maxTokens: 1024;
    betas: #('mcp-client-2025-04-04');
    mcpServers: (Array with: server);
    addUserMessage: 'List my open issues';
    yourself.

"Send. Anthropic connects out to the MCP server when the model needs it."
response := client sendMessage: params.

"The response content may carry new block types; dispatch on the predicates."
response content do: [ :block |
    block isMCPToolUseBlock ifTrue: [
        "block serverName, block name, block input carry the invocation" ].
    block isMCPToolResultBlock ifTrue: [
        "block toolUseId links back; block content is String or
         OrderedCollection of ClaudeTextBlock" ] ]
\end{lstlisting}

For the client-defined tool sibling see \S\ref{sec:tool-use}; for
Anthropic-hosted tools where the API executes the tool directly
see \S\ref{sec:server-tools}; for the opt-in header mechanism see
\S\ref{sec:beta-headers}.

\medskip
\begin{sloppypar}
\noindent{\small\textsf{Status: \emph{Available}.\quad
Anthropic docs: \url{https://docs.claude.com/en/docs/agents-and-tools/mcp-connector}.}}
\end{sloppypar}

\section{Agent Skills}
\label{sec:agent-skills}

An Agent Skill is a portable bundle of instructions and
supporting scripts that Anthropic loads into a code-execution
container on Claude's behalf during a Messages request.  Each
skill is a directory rooted at a \texttt{SKILL.md}, uploaded as a
versioned artifact through the Skills API; a request that
references the skill id inside its container spec gets the files
unpacked into the container's filesystem and made visible to the
model.  Anthropic ships a pre-built catalog --- \texttt{pptx},
\texttt{xlsx}, \texttt{docx}, \texttt{pdf} --- for document
generation, and customers can publish their own skills for any
repeatable capability they want Claude to reach for on the next
turn.  The pre-built catalog is server-side: the SDK does not
enumerate the pre-built names locally, it just loads any skill
by id into the request's container.

Skills sit alongside \S\ref{sec:tool-use} and
\S\ref{sec:server-tools} as a third kind of capability surface, and
the choice between them comes down to what the capability \emph{is}.
A \textbf{tool} is a single-purpose callable the model invokes
mid-turn through the tool-use protocol --- the model emits a
\texttt{tool\_use} block naming the tool and its arguments, the
caller executes the tool, and the model receives a
\texttt{tool\_result} on the next turn.  A \textbf{skill} is a
versioned bundle of prose and supporting scripts that Anthropic
unpacks into the request's code-execution container; the model loads
the bundle's \texttt{SKILL.md} and uses the included files however
the task demands, with no per-step round-trip back to the caller.
\textbf{Retrieval-augmented generation} (RAG) is a third neighbour
that supplies reference material the model reads but does not
execute --- a similarity search over a corpus surfaces the documents
most relevant to the current turn, and those documents land in the
context window as plain text.  Tools execute, skills are loaded and
used, RAG retrieves and reads.

The decision between the three follows the shape of the request.
``Query the weather API for Berlin'' is a tool: a single call with
typed arguments and a typed return value.  ``Generate a PowerPoint
from these bullet points'' is a skill --- Anthropic's pre-built
\texttt{pptx} catalog entry (see §\ref{sec:skills} for the catalog
shape) bundles the templates and rendering scripts the model needs,
and the call to the model carries no per-slide tool round-trips.
``Find the three most similar past incidents in our postmortem
archive'' is RAG: the search runs against a vector index, the top-k
documents enter the context, and the model reasons over them.
``Compile a quarterly sales report by region'' is again a skill ---
the \texttt{pptx} and \texttt{xlsx} catalog entries compose, the
model loads both into one container, and the output is a workbook
plus a slide deck the model assembles in one turn.

Skills, tools, and RAG compose freely.  A skill's container can
call tools the request also declares --- the model can reach for the
weather tool from inside a script the \texttt{pptx} skill loaded.
A tool the caller authors can synthesize skill output --- the tool's
implementation can post-process a workbook the model just built.
And RAG complements both: the retrieved documents land in the prompt
that drives the skill's container, or supply the context a tool's
arguments encode.  The three layers stack rather than compete.

This section is the feature-guide entry point for skills; the
mechanics --- upload, list, version, delete, container loading,
page-token pagination, multipart streaming, skill specs in
\texttt{ClaudeContainerParams} and \texttt{ClaudeSkillSpec} ---
are specified in \S\ref{sec:skills}.  This section is reference
material; no runnable example applies.  A complete skill-upload
example lives in \S\ref{sec:skills}.

For the Skills API mechanics see \S\ref{sec:skills}; for tools
as an alternate capability shape see \S\ref{sec:tool-use}.

\medskip
\begin{sloppypar}
\noindent{\small\textsf{Status: \emph{Available}.\quad
Anthropic docs: \url{https://docs.claude.com/en/docs/agents-and-tools/agent-skills/overview}.}}
\end{sloppypar}

\end{document}
